\documentclass[journal]{IEEEtran}

% *** CITATION PACKAGES ***
\usepackage[spanish]{babel}
\usepackage{graphicx}
\usepackage{float}
\usepackage{placeins}
\usepackage{amsmath,amssymb}
\usepackage{siunitx}
\usepackage{tabularx}
\usepackage{booktabs}
\usepackage{array}
\usepackage{multirow}
\usepackage{longtable}
\usepackage{listings}
\usepackage{color}
\usepackage{url}
\usepackage{dblfloatfix}

\usepackage{tikz}
\usetikzlibrary{arrows.meta,positioning,shapes.geometric,fit,calc,backgrounds}

\usepackage{caption}
\captionsetup[figure]{justification=centering}

\raggedbottom
\setlength{\textfloatsep}{8pt}
\setlength{\floatsep}{6pt}
\setlength{\intextsep}{6pt}

% Configuración para fragmentos de código.
\definecolor{codegray}{rgb}{0.45,0.45,0.45}
\definecolor{codeblue}{rgb}{0.10,0.20,0.70}
\lstset{
  basicstyle=\ttfamily\footnotesize,
  keywordstyle=\color{codeblue},
  commentstyle=\color{codegray},
  numbers=left,
  numberstyle=\tiny\color{codegray},
  stepnumber=1,
  numbersep=6pt,
  frame=single,
  breaklines=true,
  showstringspaces=false,
  tabsize=4,
  captionpos=b,
  literate=
    {á}{{\'a}}1 {é}{{\'e}}1 {í}{{\'i}}1 {ó}{{\'o}}1 {ú}{{\'u}}1
    {Á}{{\'A}}1 {É}{{\'E}}1 {Í}{{\'I}}1 {Ó}{{\'O}}1 {Ú}{{\'U}}1
    {ñ}{{\~n}}1 {Ñ}{{\~N}}1 {ü}{{\"u}}1 {Ü}{{\"U}}1
}

\title{Implementación escalar y vectorial con AVX2}
\author{Nombre del autor}

\begin{document}

\maketitle

\begin{abstract}
Este proyecto implementa un normalizador estadístico y extractor de
descriptores (suma, media, varianza, mínimo, máximo y \textit{z-score}) para
arreglos \texttt{float32} en x86-64 mediante NASM, enlazado con un controlador
de evaluación escrito en C. Se compara una implementación escalar, basada en
instrucciones SSE, con una implementación vectorial optimizada mediante AVX2 y
registros YMM de 256 bits. La evaluación cuantitativa demuestra aceleraciones
de hasta $5.80\times$ y reducciones del número de instrucciones de hasta
$7.19\times$ cuando los datos residen en los niveles de caché L1, L2 o L3. Al
trabajar desde DRAM, la aceleración converge a $2.44\times$ debido a la
saturación del ancho de banda de memoria, fenómeno conocido como
\textit{Memory Wall}.
\end{abstract}

\begin{IEEEkeywords}
SIMD, AVX2, x86-64, NASM, System V AMD64 ABI, normalización \textit{z-score},
speedup, IPC, \textit{Memory Wall}.
\end{IEEEkeywords}

\section{Introducción y objetivos}
El cálculo de estadísticos descriptivos y la normalización tipificada,
definida para cada elemento como
\begin{equation}
  z_i = \frac{x_i-\mu}{\sigma},
  \label{eq:zscore}
\end{equation}
constituyen operaciones centrales en el procesamiento de señales, la visión
artificial y el aprendizaje profundo. Ante el agotamiento del escalado de
Dennard y los límites térmicos que restringen el aumento de la frecuencia, el
incremento del rendimiento microarquitectónico depende críticamente del
paralelismo a nivel de datos (DLP, \textit{Data-Level Parallelism}). La
vectorización manual mediante extensiones SIMD, en particular AVX2, permite
procesar múltiples carriles simultáneamente con cada instrucción, reduciendo la
latencia computacional y la carga sobre el decodificador frente al modelo
secuencial escalar clásico (SISD).

\subsection{Objetivo general}
Diseñar, implementar en ensamblador x86-64 mediante NASM y evaluar
cuantitativamente dos núcleos de cómputo funcionalmente equivalentes: uno
escalar, basado en SSE, y otro vectorial, basado en AVX2. Ambos se integran con
un controlador de pruebas en C bajo la convención System V AMD64 ABI para el
cálculo de descriptores estadísticos y la normalización \textit{z-score} de
arreglos \texttt{float32}. Además, se analiza su rendimiento, aceleración e IPC
en relación con la jerarquía de memoria.

\subsection{Objetivos específicos}
\begin{enumerate}
  \item[OA1.] Explicar y contrastar el modelo de ejecución SIMD frente a SISD,
  analizando el procesamiento simultáneo de ocho elementos por registro YMM
  frente al procesamiento escalar elemento a elemento.

  \item[OA2.] Describir la evolución de las extensiones vectoriales x86
  (MMX, SSE, AVX, AVX2 y AVX-512) y fundamentar la selección de AVX2 como un
  compromiso adecuado entre compatibilidad y rendimiento.

  \item[OA3.] Desarrollar núcleos de cómputo en NASM para x86-64 mediante
  instrucciones empaquetadas, entre ellas \texttt{vmovaps}, \texttt{vaddps},
  \texttt{vsubps}, \texttt{vmulps}, \texttt{vdivps}, \texttt{vminps} y
  \texttt{vmaxps}, cumpliendo la convención System V AMD64 ABI.

  \item[OA4.] Aplicar alineación estricta de memoria a 32 bytes mediante
  \texttt{aligned\_alloc} y resolver el procesamiento del remanente mediante el
  enmascaramiento bit a bit \texttt{n \& \textasciitilde 7}, seguido de un
  bucle escalar de cierre (\textit{tail loop}).

  \item[OA5.] Emplear GDB para inspeccionar registros vectoriales con
  \texttt{print \$ymm0.v8\_float}, punteros y memoria con \texttt{x/8fw}, a fin
  de validar la corrección matemática y la ausencia de excepciones de protección
  general (\#GP).

  \item[OA6.] Medir experimentalmente la latencia mediante
  \texttt{clock\_gettime}, así como el speedup y el IPC mediante
  \texttt{perf stat}, para tamaños
  $N\in[10^3,\,2\times10^7]$. Las variaciones del factor de aceleración se
  correlacionan con la jerarquía de cachés y el cuello de botella de DRAM
  (\textit{Memory Wall}).

  \item[OA7.] Argumentar las decisiones de diseño de bajo nivel, incluido el
  árbol de reducción horizontal con \texttt{vextractf128} y la mitigación de
  las penalizaciones de transición AVX--SSE mediante \texttt{vzeroupper}.
\end{enumerate}

\subsection{Comparación conceptual entre SIMD y SISD}

La Tabla~\ref{tab:simd-sisd} resume las diferencias principales entre el
procesamiento escalar y el procesamiento vectorial empleado en este trabajo.

\begin{table*}[!t]
  \centering
  \caption{Comparación conceptual entre los modelos escalar SISD y vectorial SIMD.}
  \label{tab:simd-sisd}
  \begin{tabularx}{\textwidth}{@{}>{\raggedright\arraybackslash}p{0.20\textwidth}
    >{\raggedright\arraybackslash}X >{\raggedright\arraybackslash}X@{}}
    \toprule
    Criterio & Modelo escalar (SISD) & Modelo vectorial AVX2 (SIMD) \\
    \midrule
    Formato del dato & Escalar simple: un elemento \texttt{float32} &
      Empaquetado: ocho elementos \texttt{float32} \\
    Registro utilizado & \texttt{xmm0}--\texttt{xmm7}, usando el carril bajo
      de 32 bits & \texttt{ymm0}--\texttt{ymm5}, usando los 256 bits \\
    Operaciones por instrucción & Una operación de punto flotante &
      Ocho operaciones simultáneas por instrucción \\
    Sobrecarga de control & $N$ iteraciones con saltos y ramas &
      $\lfloor N/8\rfloor$ iteraciones vectoriales más el remanente escalar \\
    Presión sobre el decodificador & Alta, debido a la mayor emisión de
      microoperaciones & Reducida; se observaron hasta $7.19\times$ menos
      instrucciones \\
    Reducción de sumas & Inmediata mediante un acumulador secuencial &
      Árbol horizontal mediante \texttt{vextractf128} y permutaciones \\
    Alineación de memoria & Estándar, sin requisitos críticos para las cargas
      escalares & Alineación de 32 bytes para \texttt{vmovaps}, necesaria para
      evitar excepciones \#GP \\
    \bottomrule
  \end{tabularx}
\end{table*}

\subsection{Factor de aceleración teórico}

En una arquitectura SIMD que ejecuta operaciones homogéneas sobre datos
contiguos, la aceleración ideal está determinada por el cociente entre el ancho
del registro vectorial, $W_{\mathrm{vector}}$, y el tamaño del dato escalar,
$W_{\mathrm{dato}}$:
\begin{equation}
  S_{\text{teórico}}
  = \frac{T_{\mathrm{SISD}}}{T_{\mathrm{SIMD}}}
  = \frac{W_{\mathrm{vector}}}{W_{\mathrm{dato}}}
  = \frac{256\ \text{bits}}{32\ \text{bits}}
  = 8.0\times.
  \label{eq:speedup-teorico}
\end{equation}
En este caso, $W_{\mathrm{vector}}=256$ bits corresponde al ancho de un registro
YMM de AVX2, mientras que $W_{\mathrm{dato}}=32$ bits corresponde al formato
\texttt{float32} de IEEE~754. Por tanto, una sola instrucción empaquetada AVX2
puede realizar el equivalente aritmético de ocho instrucciones escalares
independientes, de modo que idealmente
\begin{equation}
  \mathrm{Throughput}_{\mathrm{SIMD}}
  = 8\,\mathrm{Throughput}_{\mathrm{SISD}}.
  \label{eq:throughput-teorico}
\end{equation}
Este límite supone ausencia de dependencias, sobrecostes de reducción y
restricciones de memoria; por ello, el speedup observado en la práctica puede
ser menor.

\section{Entorno de pruebas}

\subsection{Hardware y soporte de AVX2}

La evaluación experimental se realizó sobre un procesador Intel Core
i3-1215U de 12.\textsuperscript{a} generación, basado en la arquitectura Alder
Lake y fabricado mediante el nodo Intel~7. Este procesador presenta una
microarquitectura híbrida asimétrica compuesta por dos núcleos de rendimiento
(P-cores Golden Cove con soporte para \textit{Hyper-Threading}) y cuatro
núcleos de eficiencia (E-cores Gracemont sin multihilo), para un total de seis
núcleos físicos y ocho hilos lógicos.

La jerarquía de memoria dispone de cachés privadas L1 de datos (L1d) de
\SI{48}{\kibi\byte} por P-core y \SI{32}{\kibi\byte} por E-core. Cada P-core
cuenta con una caché L2 privada de \SI{1.25}{\mebi\byte}, mientras que el
clúster de E-cores comparte una caché L2 de \SI{2}{\mebi\byte}. El último nivel
corresponde a una caché L3 unificada de \SI{10}{\mebi\byte}, denominada Intel
Smart Cache, con líneas de caché de 64 bytes.

La inspección mediante \texttt{lscpu} y el análisis de
\texttt{/proc/cpuinfo} confirmaron el soporte nativo de AVX2 sobre registros
YMM de 256 bits, además de las extensiones FMA3 y SSE4.2. Las características
principales del sistema se resumen en la Tabla~\ref{tab:hardware}.

\begin{table}[H]
  \centering
  \caption{Características del entorno de hardware.}
  \label{tab:hardware}
  \begin{tabularx}{\columnwidth}{@{}>{\raggedright\arraybackslash}p{0.31\columnwidth}
    >{\raggedright\arraybackslash}X@{}}
    \toprule
    Característica & Valor \\
    \midrule
    Modelo de CPU & Intel Core i3-1215U (12.\textsuperscript{a} generación,
      Alder Lake, Intel~7) \\
    Arquitectura & x86-64 híbrida: 2 P-cores Golden Cove y 4 E-cores
      Gracemont \\
    Núcleos/hilos & 6 núcleos físicos (2P + 4E) y 8 hilos lógicos \\
    Frecuencia base/turbo & P-cores: \SI{1.20}{\giga\hertz}/
      \SI{4.40}{\giga\hertz}; E-cores: \SI{0.90}{\giga\hertz}/
      \SI{3.30}{\giga\hertz} \\
    Caché L1d & \SI{48}{\kibi\byte}/P-core (12 vías),
      \SI{32}{\kibi\byte}/E-core (8 vías); total:
      \SI{224}{\kibi\byte} \\
    Caché L1i & \SI{32}{\kibi\byte}/P-core (8 vías),
      \SI{64}{\kibi\byte}/E-core (8 vías); total:
      \SI{320}{\kibi\byte} \\
    Caché L2 & \SI{1.25}{\mebi\byte}/P-core y \SI{2}{\mebi\byte}
      compartidos por el clúster de E-cores; total: \SI{4.5}{\mebi\byte} \\
    Caché L3 (LLC) & \SI{10}{\mebi\byte} unificada y compartida
      (Intel Smart Cache) \\
    Línea de caché & 64 bytes \\
    Flags relevantes & \texttt{avx}, \texttt{avx2}, \texttt{fma},
      \texttt{sse4\_1}, \texttt{sse4\_2}, \texttt{bmi1}, \texttt{bmi2},
      \texttt{popcnt} \\
    \bottomrule
  \end{tabularx}
\end{table}

\subsection{Herramientas utilizadas}

El desarrollo, la instrumentación y el análisis de los núcleos de cómputo se
llevaron a cabo mediante el conjunto de herramientas nativo de GNU/Linux
x86-64. El flujo de trabajo integró el ensamblado de bajo nivel con metadatos
DWARF, la compilación modular del arnés de pruebas conforme a la convención
System V AMD64 ABI, la inspección interactiva de registros vectoriales y la
caracterización microarquitectónica no invasiva mediante la unidad de monitoreo
de rendimiento de la CPU (PMU, \textit{Performance Monitoring Unit}). La
Tabla~\ref{tab:herramientas} detalla las versiones y el propósito de cada
herramienta.

\begin{table*}[!t]
  \centering
  \caption{Herramientas utilizadas en el entorno de pruebas.}
  \label{tab:herramientas}
  \begin{tabularx}{\textwidth}{@{}>{\raggedright\arraybackslash}p{0.11\textwidth}
    >{\raggedright\arraybackslash}p{0.13\textwidth}
    >{\raggedright\arraybackslash}X@{}}
    \toprule
    Herramienta & Versión & Propósito en el proyecto \\
    \midrule
    NASM & 3.01 (compatible con 2.16+) & Ensamblado de las rutinas escalares y
      vectoriales (\texttt{stats\_scalar.asm} y \texttt{stats\_vector.asm}) en
      objetos ELF64 con símbolos de depuración DWARF mediante
      \texttt{-f elf64 -g -F dwarf}. \\
    GCC & 15.2.0 & Compilación del controlador \texttt{driver.c} con
      \texttt{-O2 -std=gnu11}, gestión de búferes alineados a 32 bytes mediante
      \texttt{aligned\_alloc} y enlace final del ejecutable. \\
    GDB & 17.1 & Depuración de bajo nivel, validación de la alineación de
      memoria mediante \texttt{addr \% 32 == 0} e inspección de los ocho
      carriles flotantes de registros YMM con
      \texttt{\$ymm0.v8\_float}. \\
    perf & 6.13 & Perfilado mediante los contadores de hardware de la CPU:
      ciclos de reloj, instrucciones retiradas, IPC, fallos de caché L1d y
      fallos en la LLC. \\
    GNU Make & 4.4.1 & Automatización del ciclo de compilación, construcción de
      los binarios \texttt{norm\_scalar} y \texttt{norm\_vector}, ejecución de
      las suites de prueba y realización sistemática de benchmarks. \\
    \bottomrule
  \end{tabularx}
\end{table*}

\section{Implementación escalar}

\subsection{Descripción del algoritmo}

El procesamiento estadístico escalar se divide en dos rutinas desacopladas,
\texttt{compute\_stats} y \texttt{normalize\_array}, implementadas íntegramente
en ensamblador x86-64. La primera calcula los descriptores estadísticos mediante
dos recorridos del arreglo, mientras que la segunda produce el arreglo
normalizado a partir de la media y la desviación estándar obtenidas.

\subsubsection{Algoritmo de dos pasadas en \texttt{compute\_stats}}

En la primera pasada se fusionan la reducción de suma y el seguimiento de los
extremos. Los tres acumuladores se inicializan con el primer elemento:
\begin{equation}
  S=x_0, \qquad x_{\min}=x_0, \qquad x_{\max}=x_0.
  \label{eq:scalar-init}
\end{equation}
Luego se recorre el arreglo desde $i=1$ hasta $N-1$, actualizando
simultáneamente
\begin{equation}
  S=\sum_{i=0}^{N-1}x_i, \qquad
  x_{\min}=\min_i(x_i), \qquad
  x_{\max}=\max_i(x_i).
  \label{eq:scalar-pass1}
\end{equation}
Al finalizar, $N$ se convierte a punto flotante y se calcula la media
aritmética:
\begin{equation}
  \mu=\frac{S}{N}=\frac{1}{N}\sum_{i=0}^{N-1}x_i.
  \label{eq:scalar-mean}
\end{equation}
La fusión de estas tres operaciones aprovecha las líneas de caché que ya se
encuentran en L1d y evita recorridos adicionales para obtener el mínimo y el
máximo.

La segunda pasada calcula la varianza poblacional. Con la media retenida en un
registro XMM se acumula la suma de desviaciones cuadráticas
\begin{equation}
  SS=\sum_{i=0}^{N-1}(x_i-\mu)^2,
  \label{eq:scalar-ss}
\end{equation}
y posteriormente se divide entre $N$:
\begin{equation}
  \sigma^2=\frac{SS}{N}
  =\frac{1}{N}\sum_{i=0}^{N-1}(x_i-\mu)^2.
  \label{eq:scalar-variance}
\end{equation}
Se descartó deliberadamente la expresión de una sola pasada
\begin{equation}
  \sigma^2=\frac{1}{N}\sum_{i=0}^{N-1}x_i^2-\mu^2,
  \label{eq:naive-variance}
\end{equation}
porque puede sufrir cancelación catastrófica al restar magnitudes grandes y muy
próximas. En precisión simple IEEE~754, cuya mantisa efectiva es de 24 bits,
esta pérdida de significancia puede amplificar el error de redondeo e incluso
producir varianzas negativas espurias.

\subsubsection{Normalización en \texttt{normalize\_array}}

La rutina recorre el arreglo una vez y aplica la transformación
\begin{equation}
  z_i=\frac{x_i-\mu}{\sigma}.
  \label{eq:scalar-normalization}
\end{equation}
Cuando $\sigma=0.0$, todos los elementos presentan dispersión nula. En este
caso, la rutina copia cada $x_i$ a la salida sin efectuar la división, con lo
que evita resultados indeterminados, valores NaN o infinitos.

\subsubsection{Complejidad algorítmica}

Ambas rutinas presentan complejidad temporal $O(N)$. En particular,
\texttt{compute\_stats} realiza dos recorridos, equivalentes a $2N$ iteraciones,
y \texttt{normalize\_array} realiza un recorrido adicional de $N$ iteraciones.
La complejidad espacial auxiliar es $O(1)$, pues el cálculo no reserva memoria
dinámica y mantiene sus acumuladores en registros de la CPU.

\subsection{Decisiones de implementación}

\subsubsection{Uso de registros escalares XMM}

La implementación emplea exclusivamente el carril inferior de 32 bits de los
registros \texttt{xmm0}--\texttt{xmm6}. Las instrucciones con sufijo
\texttt{ss} (\textit{Scalar Single-Precision}), como \texttt{addss},
\texttt{subss}, \texttt{mulss} y \texttt{divss}, procesan un elemento
\texttt{float32} por instrucción y corresponden al modelo SISD utilizado como
línea base.

\subsubsection{Convención System V AMD64 ABI}

En \texttt{compute\_stats}, los parámetros se reciben en el orden establecido
por la ABI: \texttt{rdi} contiene \texttt{arr}, \texttt{rsi} contiene $N$ y
\texttt{rdx}, \texttt{rcx}, \texttt{r8} y \texttt{r9} contienen los punteros
de salida para la media, la varianza, el mínimo y el máximo, respectivamente.
En \texttt{normalize\_array}, \texttt{rdi}, \texttt{rsi} y \texttt{edx}
contienen los punteros de entrada y salida y la longitud; los argumentos
flotantes $\mu$ y $\sigma$ se reciben en \texttt{xmm0} y \texttt{xmm1}.

Para conservar los punteros durante las dos pasadas de
\texttt{compute\_stats}, se trasladan a \texttt{r12}--\texttt{r15},
\texttt{rbx} y \texttt{rbp}. Estos registros no volátiles
(\textit{callee-saved}) se guardan mediante \texttt{push} en el prólogo y se
restauran con \texttt{pop}, en orden inverso, antes del retorno. De esta forma
se preserva el estado esperado por la función invocadora escrita en C.

\subsubsection{Selección de extremos sin bifurcaciones}

Las instrucciones \texttt{minss} y \texttt{maxss} actualizan los extremos
directamente en hardware, sin recurrir a \texttt{comiss} ni a saltos
condicionales como \texttt{ja} o \texttt{jb}. Este enfoque
\textit{branchless} evita penalizaciones por predicciones de salto incorrectas,
que en procesadores modernos pueden costar del orden de decenas de ciclos.

\subsubsection{Detección de varianza nula}

Antes de normalizar, \texttt{normalize\_array} construye la constante $0.0$
mediante \texttt{xorps xmm2, xmm2} y la compara con la desviación estándar
usando \texttt{ucomiss xmm1, xmm2}. Si $\sigma=0.0$, el control se desvía a
\texttt{.copy\_loop}; así se evita la división por cero y la generación de NaN
o $\pm\infty$.

\subsection{Fragmentos de código comentados}
\begin{lstlisting}[language={[x86masm]Assembler},caption={Bucles principales de \texttt{stats\_scalar.asm}.},label={lst:escalar}]
; PASADA 1: suma, minimo y maximo
movss   xmm0, [r12]         ; suma = arr[0]
movaps  xmm1, xmm0          ; min = arr[0]
movaps  xmm2, xmm0          ; max = arr[0]
mov     eax, 1              ; i = 1
.pass1_loop:
    cmp     eax, r13d       ; i >= n?
    jge     .pass1_done
    movss   xmm3, [r12 + rax*4] ; arr[i]
    addss   xmm0, xmm3      ; suma += arr[i]
    minss   xmm1, xmm3      ; min = min(min, arr[i])
    maxss   xmm2, xmm3      ; max = max(max, arr[i])
    inc     eax
    jmp     .pass1_loop
.pass1_done:
    cvtsi2ss xmm4, r13d     ; xmm4 = (float)n
    movaps  xmm5, xmm0
    divss   xmm5, xmm4      ; mean = suma / n
    movss   [rbx], xmm1     ; *min = min
    movss   [rbp], xmm2     ; *max = max
    movss   [r14], xmm5     ; *mean = mean

; PASADA 2: varianza poblacional
xor     eax, eax
xorps   xmm6, xmm6          ; acumulador = 0.0
.pass2_loop:
    cmp     eax, r13d
    jge     .pass2_done
    movss   xmm3, [r12 + rax*4]
    subss   xmm3, xmm5      ; arr[i] - mean
    mulss   xmm3, xmm3      ; (arr[i] - mean)^2
    addss   xmm6, xmm3
    inc     eax
    jmp     .pass2_loop
.pass2_done:
    divss   xmm6, xmm4      ; var = SS / n
    movss   [r15], xmm6     ; *var = var

; NORMALIZACION: z-score o copia si sigma = 0
; rdi=in, rsi=out, edx=n, xmm0=mean, xmm1=stddev
test    edx, edx
jle     .norm_done
xorps   xmm2, xmm2
ucomiss xmm1, xmm2
je      .copy_loop
xor     eax, eax
.norm_loop:
    cmp     eax, edx
    jge     .norm_done
    movss   xmm3, [rdi + rax*4]
    subss   xmm3, xmm0
    divss   xmm3, xmm1
    movss   [rsi + rax*4], xmm3
    inc     eax
    jmp     .norm_loop
.copy_loop:
    xor     eax, eax
.copy_inner:
    cmp     eax, edx
    jge     .norm_done
    movss   xmm3, [rdi + rax*4]
    movss   [rsi + rax*4], xmm3
    inc     eax
    jmp     .copy_inner
.norm_done:
    ret
\end{lstlisting}

La Tabla~\ref{tab:registros-escalar} resume la asignación de registros en las
dos rutinas. Para conservarla dentro de una sola columna se utiliza una
tipografía compacta y descripciones abreviadas.

\begin{table}[H]
  \centering
  \caption{Asignación de registros en la implementación escalar.}
  \label{tab:registros-escalar}
  \scriptsize
  \setlength{\tabcolsep}{2.5pt}
  \begin{tabularx}{\columnwidth}{@{}>{\raggedright\arraybackslash}p{0.13\columnwidth}
    >{\raggedright\arraybackslash}p{0.16\columnwidth}
    >{\raggedright\arraybackslash}X
    >{\raggedright\arraybackslash}X@{}}
    \toprule
    Registro & Tipo & \texttt{compute\_stats} & \texttt{normalize\_array} \\
    \midrule
    \texttt{r12}/\texttt{rdi} & Entero 64 b & Base de \texttt{arr} & Base de
      \texttt{in} \\
    \texttt{rsi} & Entero 64 b & No usado & Base de \texttt{out} \\
    \texttt{r13d}/\texttt{edx} & Entero 32 b & Longitud $N$ & Longitud $N$ \\
    \texttt{eax}/\texttt{rax} & Entero 32/64 b & Índice $i$ & Índice $i$ \\
    \texttt{xmm0} & Escalar 32 b & Suma $\sum x_i$ & Media $\mu$ \\
    \texttt{xmm1} & Escalar 32 b & Mínimo & Desviación $\sigma$ \\
    \texttt{xmm2} & Escalar 32 b & Máximo & Constante $0.0$ \\
    \texttt{xmm3} & Escalar 32 b & Operando temporal & Operando temporal \\
    \texttt{xmm4} & Escalar 32 b & $N$ convertido a float & No usado \\
    \texttt{xmm5} & Escalar 32 b & Media $\mu$ & No usado \\
    \texttt{xmm6} & Escalar 32 b & $SS$ y varianza $\sigma^2$ & No usado \\
    \texttt{r14}, \texttt{r15},
      \texttt{rbx}, \texttt{rbp} & Punteros 64 b & Destinos de media,
      varianza, mínimo y máximo & No usados \\
    \bottomrule
  \end{tabularx}
\end{table}

\subsection{Evidencia visual pendiente}

Los siguientes espacios quedan reservados para incorporar las capturas sin
alterar posteriormente la estructura o numeración del informe.

\begin{figure}[H]
  \centering
  % Reemplace el recuadro por:
  % \includegraphics[width=\columnwidth]{figuras/gdb-registros-xmm.png}
  \fbox{\parbox[c][3.5cm][c]{0.92\columnwidth}{\centering
    Captura pendiente: sesión de GDB inspeccionando los registros
    \texttt{xmm0}--\texttt{xmm6} y la memoria de salida después de ejecutar
    \texttt{compute\_stats}.}}
  \caption{Inspección de registros XMM y memoria de salida mediante GDB.}
  \label{fig:gdb-escalar-xmm}
\end{figure}

\begin{figure}[H]
  \centering
  % Reemplace el recuadro por:
  % \includegraphics[width=\columnwidth]{figuras/salida-norm-scalar-n16.png}
  \fbox{\parbox[c][3.2cm][c]{0.92\columnwidth}{\centering
    Captura pendiente: salida de \texttt{./bin/norm\_scalar} para el caso
    $N=16$, con validación de las estadísticas calculadas.}}
  \caption{Ejecución y validación de la versión escalar para $N=16$.}
  \label{fig:salida-escalar-n16}
\end{figure}

\section{Implementación vectorial con AVX2}

\subsection{Elección de AVX2}

El conjunto de instrucciones AVX2 (\textit{Advanced Vector Extensions 2}) se
seleccionó después de evaluar la capacidad, disponibilidad y coste operativo de
las principales extensiones SIMD de x86-64.

Frente a SSE, que utiliza registros XMM de 128 bits y procesa cuatro valores
\texttt{float32} por instrucción, AVX2 duplica el ancho a 256 bits y permite
operar sobre ocho valores simultáneamente. En consecuencia, el techo ideal
asociado únicamente al ancho vectorial aumenta de $4.0\times$ a $8.0\times$.
Además, la codificación tradicional de SSE emplea instrucciones destructivas de
dos operandos, por ejemplo \texttt{addps dest, src}, mientras que el prefijo VEX
de AVX2 permite instrucciones no destructivas de tres operandos:
\begin{equation}
  \begin{aligned}
    &\texttt{vaddps dest, src1, src2} \\
    &\Longrightarrow\quad
      \texttt{dest}=\texttt{src1}+\texttt{src2}.
  \end{aligned}
  \label{eq:vex-three-operands}
\end{equation}
Esto preserva las fuentes y reduce la necesidad de copias auxiliares con
\texttt{movaps}.

AVX-512 elevaría el ancho a 512 bits, equivalente a 16 valores
\texttt{float32}. Sin embargo, no está disponible de forma homogénea en el
Intel Core i3-1215U utilizado: sus E-cores Gracemont no incorporan unidades
AVX-512 y la plataforma mantiene deshabilitada esta extensión para ofrecer un
repertorio uniforme entre núcleos. En las plataformas que sí la implementan,
las cargas AVX-512 intensivas también pueden provocar reducciones de frecuencia
dependientes del procesador y del presupuesto térmico. A esto se añade la
fragmentación de AVX-512 en subconjuntos como F, CD, BW, DQ y VL. AVX2 ofrece,
por tanto, una cobertura más amplia y predecible en equipos x86-64 modernos.

La Tabla~\ref{tab:comparacion-simd} resume los criterios de selección. Se
mantiene dentro de una columna mediante descripciones compactas.

\begin{table}[H]
  \centering
  \caption{Comparación de extensiones vectoriales x86.}
  \label{tab:comparacion-simd}
  \scriptsize
  \setlength{\tabcolsep}{2pt}
  \begin{tabularx}{\columnwidth}{@{}>{\raggedright\arraybackslash}p{0.25\columnwidth}
    >{\centering\arraybackslash}X >{\centering\arraybackslash}X
    >{\centering\arraybackslash}X@{}}
    \toprule
    Criterio & SSE & AVX2 & AVX-512 \\
    \midrule
    Registro & 128 b XMM & 256 b YMM & 512 b ZMM \\
    \texttt{float32}/vector & 4 & 8 & 16 \\
    Operandos & 2, destructivos & 3, VEX & 3, EVEX y máscara \\
    Techo ideal & $4\times$ & $8\times$ & $16\times$ \\
    Compatibilidad & Muy amplia & Amplia & Restringida en consumo \\
    Efecto térmico & Bajo & Bajo o moderado & Puede ser alto \\
    \bottomrule
  \end{tabularx}
\end{table}

En síntesis, AVX2 constituye el punto de equilibrio para este trabajo: ofrece
ocho operaciones simultáneas por instrucción, está soportado por todos los
núcleos del sistema experimental y evita depender de una extensión ausente en
la plataforma objetivo.

\subsection{Organización del procesamiento vectorial}

Los datos se agrupan en bloques de ocho números IEEE~754 de precisión simple.
Cada registro YMM se interpreta como ocho carriles independientes:
\begin{equation}
  \mathrm{YMM}=
  [f_0,f_1,f_2,f_3\mid f_4,f_5,f_6,f_7].
  \label{eq:ymm-lanes}
\end{equation}
Los acumuladores principales ocupan \texttt{ymm0}--\texttt{ymm5}. Cada
iteración consume $8\times4=32$ bytes contiguos; el índice se conserva en
\texttt{eax}/\texttt{rax} y avanza mediante \texttt{add eax, 8}. El modo de
direccionamiento SIB convierte el índice de elemento en desplazamiento de bytes
con expresiones como \texttt{[rdi + rax*4]} o
\texttt{[r12 + rax*4]}.

Las principales instrucciones VEX utilizadas son las siguientes:
\begin{itemize}
  \item \texttt{vmovaps}: carga o almacena ocho valores empaquetados. Las
  direcciones usadas por esta implementación se mantienen alineadas a 32 bytes,
  condición necesaria para esta variante alineada.
  \item \texttt{vaddps}, \texttt{vsubps}, \texttt{vmulps} y
  \texttt{vdivps}: ejecutan suma, resta, multiplicación y división en los ocho
  carriles.
  \item \texttt{vminps} y \texttt{vmaxps}: actualizan en paralelo los ocho
  mínimos o máximos parciales.
  \item \texttt{vbroadcastss}: replica el valor escalar del carril inferior de
  un XMM en los ocho carriles de un YMM, es decir,
  $[s,s,s,s,s,s,s,s]$.
\end{itemize}

\begin{figure}[H]
  \centering
  % \includegraphics[width=\columnwidth]{figuras/gdb-ymm0-carriles.png}
  \fbox{\parbox[c][3.1cm][c]{0.92\columnwidth}{\centering
    Captura pendiente: inspección en GDB de los ocho carriles de
    \texttt{ymm0} mediante \texttt{print \$ymm0.v8\_float} durante el bucle
    vectorial.}}
  \caption{Inspección de los carriles de un registro YMM.}
  \label{fig:gdb-ymm-lanes}
\end{figure}

\begin{figure}[H]
  \centering
  % \includegraphics[width=\columnwidth]{figuras/gdb-alineacion-32.png}
  \fbox{\parbox[c][3.1cm][c]{0.92\columnwidth}{\centering
    Captura pendiente: volcado \texttt{x/8fw arr} y comprobación de que la
    dirección base del arreglo es múltiplo de 32 bytes.}}
  \caption{Validación en GDB de la alineación del arreglo de entrada.}
  \label{fig:gdb-alineacion-avx2}
\end{figure}

\subsection{Reducción horizontal}

Al concluir el bucle vectorial, cada acumulador contiene ocho resultados
parciales:
\begin{equation}
  \mathrm{YMM}=[c_0,c_1,c_2,c_3\mid c_4,c_5,c_6,c_7].
  \label{eq:partial-lanes}
\end{equation}
Para producir un único resultado escalar se aplican árboles de reducción de
profundidad logarítmica.

\subsubsection{Reducción de suma}

Primero, \texttt{vextractf128} extrae los carriles 4--7. La instrucción
\texttt{vaddps} combina ambas mitades y deja cuatro sumas parciales:
\begin{equation}
  [c_0+c_4,\;c_1+c_5,\;c_2+c_6,\;c_3+c_7].
  \label{eq:sum-halves}
\end{equation}
Dos aplicaciones de \texttt{vhaddps} colapsan primero los pares adyacentes y
luego las dos sumas restantes, por lo que
\begin{equation}
  \texttt{xmm0}[0]=\sum_{i=0}^{7}c_i.
  \label{eq:horizontal-sum}
\end{equation}

\subsubsection{Reducción de mínimo y máximo}

x86-64 no proporciona instrucciones \texttt{vhminps} o \texttt{vhmaxps}. La
reducción se construye mediante \texttt{vextractf128}, operaciones de mínimo o
máximo y dos permutaciones. La máscara \texttt{0x4E} reorganiza los carriles
como $[2,3,0,1]$, mientras que \texttt{0xB1} los reorganiza como
$[1,0,3,2]$. El siguiente fragmento muestra la reducción del mínimo:

\begin{lstlisting}[language={[x86masm]Assembler},caption={Árbol horizontal para reducir el mínimo.},label={lst:reduce-min}]
vextractf128 xmm3, ymm1, 1
vminps       xmm1, xmm1, xmm3
vshufps      xmm3, xmm1, xmm1, 0x4E
vminps       xmm1, xmm1, xmm3
vshufps      xmm3, xmm1, xmm1, 0xB1
vminps       xmm1, xmm1, xmm3
; xmm1[0] contiene el minimo global
\end{lstlisting}

La reducción del máximo repite el mismo árbol sobre \texttt{ymm2}/
\texttt{xmm2}, sustituyendo \texttt{vminps} por \texttt{vmaxps}.

\begin{figure}[H]
  \centering
  % \includegraphics[width=\columnwidth]{figuras/gdb-reduccion-minimo.png}
  \fbox{\parbox[c][3.1cm][c]{0.92\columnwidth}{\centering
    Captura pendiente: contenido de \texttt{xmm1} después de la secuencia
    \texttt{vshufps}/\texttt{vminps}, mostrando el mínimo en el carril 0.}}
  \caption{Validación de la reducción horizontal del mínimo.}
  \label{fig:gdb-reduccion-min}
\end{figure}

\subsection{Manejo del remanente}

Cuando $N$ no es múltiplo de ocho, una carga vectorial final accedería a
posiciones ajenas al arreglo y podría atravesar un límite de página. Para
evitarlo, se calcula primero el número de elementos que forman bloques
completos:
\begin{equation}
  N_{\mathrm{vec}}=N\mathbin{\&}(\mathord{\sim}7)
  =N-(N\bmod 8).
  \label{eq:n-vector-limit}
\end{equation}
La máscara de 32 bits es
\begin{equation}
  \mathord{\sim}7=\texttt{0xFFFFFFF8},
  \label{eq:tail-mask}
\end{equation}
de modo que \texttt{and ecx, \textasciitilde 7} apaga los tres bits menos
significativos y redondea $N$ hacia abajo al múltiplo de ocho más cercano.

Si $N<8$, entonces $N_{\mathrm{vec}}=0$ y el bucle AVX2 se omite por completo.
Para $N\bmod8\in\{1,\ldots,7\}$, un bucle escalar de cierre comienza en
$i=N_{\mathrm{vec}}$ y procesa los elementos restantes con instrucciones VEX
escalares: \texttt{vmovss}, \texttt{vaddss}, \texttt{vminss},
\texttt{vmaxss}, \texttt{vsubss}, \texttt{vmulss} y \texttt{vdivss}. Cada
\texttt{vmovss} lee exactamente cuatro bytes y los resultados se acumulan sobre
los registros XMM que ya contienen la reducción horizontal.

\begin{figure}[H]
  \centering
  % \includegraphics[width=\columnwidth]{figuras/gdb-tail-loop.png}
  \fbox{\parbox[c][3.1cm][c]{0.92\columnwidth}{\centering
    Captura pendiente: depuración del \textit{tail loop} mostrando el avance
    elemento a elemento con \texttt{vmovss}.}}
  \caption{Procesamiento escalar del remanente en GDB.}
  \label{fig:gdb-tail-loop}
\end{figure}

\subsection{Fragmentos de código comentados}

Los fragmentos siguientes corresponden a los núcleos de
\texttt{asm/vector/stats\_vector.asm}.

\begin{lstlisting}[language={[x86masm]Assembler},caption={Primera pasada vectorial y reducción de suma y mínimo.},label={lst:avx2-pass1}]
.pass1_vec_loop:
    cmp     eax, ecx
    jge     .pass1_vec_reduce
    vmovaps ymm3, [r12 + rax*4]
    vaddps  ymm0, ymm0, ymm3
    vminps  ymm1, ymm1, ymm3
    vmaxps  ymm2, ymm2, ymm3
    add     eax, 8
    jmp     .pass1_vec_loop

.pass1_vec_reduce:
    ; Suma: ymm0 -> xmm0[0]
    vextractf128 xmm3, ymm0, 1
    vaddps  xmm0, xmm0, xmm3
    vhaddps xmm0, xmm0, xmm0
    vhaddps xmm0, xmm0, xmm0

    ; Minimo: ymm1 -> xmm1[0]
    vextractf128 xmm3, ymm1, 1
    vminps  xmm1, xmm1, xmm3
    vshufps xmm3, xmm1, xmm1, 0x4E
    vminps  xmm1, xmm1, xmm3
    vshufps xmm3, xmm1, xmm1, 0xB1
    vminps  xmm1, xmm1, xmm3
\end{lstlisting}

\begin{lstlisting}[language={[x86masm]Assembler},caption={Bucle escalar de cierre.},label={lst:avx2-tail}]
.pass1_tail_loop:
    cmp     eax, r13d
    jge     .pass1_done
    vmovss  xmm3, [r12 + rax*4]
    vaddss  xmm0, xmm0, xmm3
    vminss  xmm1, xmm1, xmm3
    vmaxss  xmm2, xmm2, xmm3
    inc     eax
    jmp     .pass1_tail_loop
\end{lstlisting}

\begin{lstlisting}[language={[x86masm]Assembler},caption={Normalización vectorial con difusión de escalares.},label={lst:avx2-normalize}]
normalize_array:
    test    edx, edx
    jle     .norm_vec_ret
    vxorps   xmm2, xmm2, xmm2
    vucomiss xmm1, xmm2
    je      .norm_copy_loop

    vbroadcastss ymm4, xmm0 ; media
    vbroadcastss ymm5, xmm1 ; desviacion
    mov     ecx, edx
    and     ecx, ~7
    xor     eax, eax
.norm_vec_loop:
    cmp     eax, ecx
    jge     .norm_scalar_tail
    vmovaps ymm0, [rdi + rax*4]
    vsubps  ymm0, ymm0, ymm4
    vdivps  ymm0, ymm0, ymm5
    vmovaps [rsi + rax*4], ymm0
    add     eax, 8
    jmp     .norm_vec_loop
\end{lstlisting}

\subsubsection{Uso de \texttt{vzeroupper}}

Después de ejecutar instrucciones AVX de 256 bits, los bits superiores de los
registros YMM pueden conservar datos activos. En microarquitecturas sensibles a
la transición entre AVX y SSE heredado, retornar directamente a código que usa
instrucciones SSE sin prefijo VEX puede introducir una penalización. Por ello,
todos los caminos de salida ejecutan
\begin{lstlisting}[language={[x86masm]Assembler},caption={Limpieza del estado superior de los registros YMM.},label={lst:vzeroupper}]
.norm_vec_done:
    vzeroupper
    ret
\end{lstlisting}
La instrucción marca como limpios los carriles superiores de los registros YMM
y evita depender del comportamiento concreto del código C o de las bibliotecas
invocadas después del retorno.

\begin{figure}[H]
  \centering
  % \includegraphics[width=\columnwidth]{figuras/gdb-vzeroupper.png}
  \fbox{\parbox[c][3.1cm][c]{0.92\columnwidth}{\centering
    Captura pendiente: desensamblado de \texttt{compute\_stats} o
    \texttt{normalize\_array} mostrando \texttt{vzeroupper} inmediatamente
    antes de \texttt{ret}.}}
  \caption{Ubicación de \texttt{vzeroupper} en el epílogo AVX2.}
  \label{fig:gdb-vzeroupper}
\end{figure}

\subsection{Asignación de registros vectoriales}

Todos los registros XMM/YMM son volátiles (\textit{caller-saved}) bajo System V
AMD64 ABI. Por ello, las rutinas pueden utilizarlos sin guardarlos previamente
en la pila. Para conservar las tablas a media columna, la asignación se divide
por función.

\begin{table}[H]
  \centering
  \caption{Registros de suma y cálculo estadístico AVX2.}
  \label{tab:registros-avx2-stats}
  \scriptsize
  \setlength{\tabcolsep}{2.5pt}
  \begin{tabularx}{\columnwidth}{@{}>{\raggedright\arraybackslash}p{0.18\columnwidth}
    >{\raggedright\arraybackslash}p{0.29\columnwidth}
    >{\raggedright\arraybackslash}X@{}}
    \toprule
    Registro & \texttt{sum\_array} & \texttt{compute\_stats} \\
    \midrule
    YMM0/XMM0 & Suma vectorial y retorno escalar & P1: suma; P2: suma de
      desviaciones cuadráticas \\
    YMM1/XMM1 & Carga vectorial o elemento del remanente & P1: mínimo; P2:
      carga y $(x-\mu)^2$ \\
    YMM2/XMM2 & Mitad alta durante la reducción & P1: máximo \\
    YMM3/XMM3 & No usado & Carga vectorial, auxiliar de permutación y difusión
      de $\mu$ \\
    YMM4/XMM4 & No usado & Divisor $N$ convertido a punto flotante \\
    YMM5/XMM5 & No usado & Media escalar $\mu=S/N$ \\
    \bottomrule
  \end{tabularx}
\end{table}

\begin{table}[H]
  \centering
  \caption{Registros de la normalización AVX2.}
  \label{tab:registros-avx2-normalize}
  \scriptsize
  \setlength{\tabcolsep}{3pt}
  \begin{tabularx}{\columnwidth}{@{}>{\raggedright\arraybackslash}p{0.24\columnwidth}
    >{\raggedright\arraybackslash}X@{}}
    \toprule
    Registro & Rol en \texttt{normalize\_array} \\
    \midrule
    YMM0/XMM0 & Carga, resta, división y almacenamiento de ocho valores \\
    YMM1/XMM1 & Argumento escalar $\sigma$ \\
    YMM2/XMM2 & Comparación $\sigma=0.0$ y temporal del remanente \\
    YMM3/XMM3 & Temporal escalar del remanente \\
    YMM4/XMM4 & Difusión $[\mu,\ldots,\mu]$ \\
    YMM5/XMM5 & Difusión $[\sigma,\ldots,\sigma]$ \\
    \bottomrule
  \end{tabularx}
\end{table}

\section{Casos de prueba}

\subsection{Comandos de reproducibilidad y ejecución}

La compilación desde cero y la verificación funcional automatizada se realizan
desde el directorio raíz del proyecto mediante los siguientes comandos:

\begin{lstlisting}[language=bash,caption={Compilación y ejecución de la suite completa.},label={lst:run-tests}]
# Limpiar y compilar las dos versiones
make clean && make

# Ejecutar la suite funcional completa
python3 tools/run_all_tests.py
\end{lstlisting}

Para reproducir manualmente el caso con $N=15$, que combina un bloque AVX2 de
ocho elementos con siete elementos de remanente, se ejecutan ambos binarios en
modo detallado y se comparan sus resúmenes estadísticos:

\begin{lstlisting}[language=bash,caption={Ejecución manual del caso con $N=15$.},label={lst:run-n15}]
# Kernel escalar, modo verbose
./bin/norm_scalar data/test_suite/tc-05_in.dat data/test_suite/tc-05_sc.dat 1

# Kernel vectorial AVX2, modo verbose
./bin/norm_vector data/test_suite/tc-05_in.dat data/test_suite/tc-05_vc.dat 1

# Comparar los resumenes estadisticos
cat data/test_suite/tc-05_sc.dat.stats.txt
cat data/test_suite/tc-05_vc.dat.stats.txt
\end{lstlisting}

\subsection{Criterio de validación y tolerancia numérica}

La equivalencia entre las implementaciones en ensamblador y el oráculo de
referencia escrito en Python se acepta cuando
\begin{equation}
  \mathrm{Err}_{\mathrm{rel}}\leq 1.0\times10^{-4}.
  \label{eq:test-tolerance}
\end{equation}
Para un valor producido por ensamblador, $v_{\mathrm{asm}}$, y su referencia,
$v_{\mathrm{ref}}$, la discrepancia relativa se calcula como
\begin{equation}
  \mathrm{Err}_{\mathrm{rel}}=
  \frac{|v_{\mathrm{asm}}-v_{\mathrm{ref}}|}
       {|v_{\mathrm{ref}}|+10^{-12}}.
  \label{eq:relative-error}
\end{equation}
El término $10^{-12}$ estabiliza el denominador cuando la referencia es nula o
muy cercana a cero.

\subsubsection{No asociatividad en IEEE 754}

La suma en punto flotante \texttt{binary32} no satisface, en general, la
propiedad asociativa:
\begin{equation}
  (a+b)+c\neq a+(b+c).
  \label{eq:non-associativity}
\end{equation}
El núcleo escalar acumula los elementos secuencialmente,
\begin{equation}
  S_{\mathrm{sc}}=(((x_0+x_1)+x_2)+\cdots+x_{N-1}),
  \label{eq:scalar-order}
\end{equation}
mientras que AVX2 mantiene ocho sumas parciales y las combina mediante un árbol
horizontal:
\begin{equation}
  S_{\mathrm{vec}}=
  \sum_{j=0}^{7}
  \left(\sum_{k=0}^{\lfloor N/8\rfloor-1}x_{8k+j}\right)
  +S_{\mathrm{tail}}.
  \label{eq:vector-order}
\end{equation}
Cada operación intermedia redondea el resultado a la precisión representable.
Por ello, la reasociación inherente al procesamiento SIMD puede producir
diferencias del orden de $10^{-7}$ a $10^{-6}$, sin que estas impliquen un
error algorítmico.

\subsection{Resultados de los casos de prueba}

La Tabla~\ref{tab:pruebas} resume los seis casos borde obligatorios. Debido al
número de resultados comparados, esta tabla utiliza excepcionalmente el ancho
de las dos columnas.

\begin{table*}[!t]
  \centering
  \caption{Resultados de los casos de prueba.}
  \label{tab:pruebas}
  \scriptsize
  \setlength{\tabcolsep}{3pt}
  \begin{tabularx}{\textwidth}{@{}c>{\raggedright\arraybackslash}p{0.21\textwidth}
    >{\raggedright\arraybackslash}X
    >{\raggedright\arraybackslash}X
    >{\raggedright\arraybackslash}Xc@{}}
    \toprule
    Caso & Entrada & Referencia & Escalar & AVX2 & Resultado \\
    \midrule
    1 & Arreglo vacío, $N=0$; archivo binario sin elementos &
      Retorno seguro: $\mu=0.0$, $\sigma^2=0.0$ &
      $\mu=0.0$, $\sigma^2=0.0$ &
      $\mu=0.0$, $\sigma^2=0.0$ & Pasa \\

    2 & Elemento único, $N=1$; $x_0=42.0$ &
      $\mu=42.0$, $\sigma^2=0.0$; salida $[42.0]$ &
      $\mu=42.0$, $\sigma^2=0.0$; salida $[42.0]$ &
      $\mu=42.0$, $\sigma^2=0.0$; salida $[42.0]$ & Pasa \\

    3 & Remanente puro, $N=7<8$; valores en $[-100,100]$ &
      $\mu=-0.6737$, $\sigma^2=1547.18$ &
      $\mu=-0.6737$, $\sigma^2=1547.18$ &
      $\mu=-0.6737$, $\sigma^2=1547.18$ & Pasa \\

    4 & Vector y remanente, $N=15$; un bloque AVX2 y siete elementos en el
      \textit{tail} &
      $\mu=19.4100$, $\sigma^2=2381.27$ &
      $\mu=19.4100$, $\sigma^2=2381.27$ &
      $\mu=19.4100$, $\sigma^2=2381.27$ & Pasa \\

    5 & Varianza nula, $N=1000$; $x_i=5.0$ para todo $i$ &
      $\mu=5.0$, $\sigma^2=0.0$; salida clonada &
      $\mu=5.0$, $\sigma^2=0.0$; salida clonada &
      $\mu=5.0$, $\sigma^2=0.0$; salida clonada & Pasa \\

    6 & Valores extremos, $N=70$; rango $[-10^6,10^6]$ con valores del orden
      de $\pm10^{-4}$ &
      $\mu=0.0$, $\sigma^2=2.8571\times10^{11}$ &
      $\mu=0.0$, $\sigma^2=2.8571\times10^{11}$ &
      $\mu=0.0$, $\sigma^2=2.8571\times10^{11}$ & Pasa \\
    \bottomrule
  \end{tabularx}
\end{table*}

\subsection{Análisis técnico de los casos borde}

\begin{enumerate}
  \item \textbf{Arreglo vacío ($N=0$):} las guardas del controlador y de las
  rutinas en ensamblador se activan antes de cualquier división o acceso a los
  punteros. Se retornan estadísticas nulas sin provocar SIGSEGV ni SIGFPE.

  \item \textbf{Elemento único ($N=1$):} la varianza es exactamente cero. La
  comparación con \texttt{ucomiss}/\texttt{vucomiss} activa la ruta de copia y
  conserva el elemento original sin dividir por cero.

  \item \textbf{Remanente puro ($N=7$):} como $N<8$, se obtiene
  $N\mathbin{\&}(\mathord{\sim}7)=0$. El bucle AVX2 se omite y todos los
  elementos se procesan mediante el \textit{tail loop} escalar.

  \item \textbf{Vector y remanente ($N=15$):} se procesa un bloque de ocho
  valores en registros YMM y luego los siete valores restantes de forma
  escalar, comprobando la transición entre ambas rutas.

  \item \textbf{Varianza nula ($N=1000$):} el arreglo constante comprueba que
  el cálculo produzca $\sigma^2=0.0$ y que la guarda impida ejecutar
  \texttt{vdivps} con un divisor nulo.

  \item \textbf{Valores extremos ($N=70$):} se mezclan signos y amplitudes
  desde $10^{-4}$ hasta $10^6$. El caso verifica que no aparezcan
  desbordamientos ni subdesbordamientos destructivos en los resultados
  esperados.
\end{enumerate}

\subsection{Dictamen de correctitud}

Los seis casos evaluados superaron la validación funcional y la comparación
numérica con el oráculo de Python. El mayor error relativo registrado entre las
implementaciones fue
\begin{equation}
  \mathrm{Err}_{\mathrm{rel,max}}=1.19\times10^{-6}
  \ll 1.0\times10^{-4}.
  \label{eq:max-relative-error}
\end{equation}
Este resultado se encuentra aproximadamente dos órdenes de magnitud por debajo
del umbral establecido y respalda la exactitud numérica de ambas versiones para
los casos considerados.

\begin{figure}[H]
  \centering
  % \includegraphics[width=\columnwidth]{figuras/suite-pruebas-pasa.png}
  \fbox{\parbox[c][3.3cm][c]{0.92\columnwidth}{\centering
    Captura pendiente: terminal ejecutando
    \texttt{python3 tools/run\_all\_tests.py} y mostrando los seis casos con
    resultado PASA.}}
  \caption{Ejecución automatizada de la suite de pruebas funcionales.}
  \label{fig:suite-pruebas}
\end{figure}

\section{Resultados de rendimiento}

\subsection{Metodología de medición}

La instrumentación temporal se realizó mediante la función POSIX
\texttt{clock\_gettime(CLOCK\_MONOTONIC)}, utilizando marcas de tiempo de alta
resolución. La región medida incluye exclusivamente el cómputo de los núcleos
\texttt{sum\_array}, \texttt{compute\_stats} y
\texttt{normalize\_array}; las operaciones de entrada y salida, como
\texttt{fread} y \texttt{fwrite}, se excluyeron de la ventana de medición.

Para cada tamaño se ejecutaron $R=30$ repeticiones independientes. Se reportan
la media muestral $\overline{T}$ y la desviación estándar muestral $s$, lo que
permite cuantificar la dispersión ocasionada por el sistema operativo, la
variación de frecuencia y los efectos térmicos. La ejecución se fijó al núcleo
P0 mediante \texttt{taskset -c 0}, evitando migraciones entre los P-cores Golden
Cove y los E-cores Gracemont del procesador híbrido.

Se seleccionaron cuatro tamaños para recorrer la jerarquía de memoria. La
huella indicada considera los arreglos de entrada y salida, ambos en formato
\texttt{float32}:
\begin{itemize}
  \item $N=10^3$: \SI{7.8}{\kibi\byte}, residente en L1d.
  \item $N=10^5$: \SI{781.2}{\kibi\byte}, residente en L2.
  \item $N=10^6$: \SI{7.63}{\mebi\byte}, residente en L3.
  \item $N=2\times10^7$: \SI{152.6}{\mebi\byte}, muy superior a la LLC y,
  por tanto, atendido principalmente desde DRAM.
\end{itemize}

\subsection{Tiempos de ejecución}

La aceleración experimental se define como
\begin{equation}
  S=\frac{\overline{T}_{\mathrm{sc}}}
          {\overline{T}_{\mathrm{vec}}},
  \label{eq:measured-speedup}
\end{equation}
y la eficiencia respecto al ancho ideal de ocho carriles como
\begin{equation}
  \eta_{\mathrm{SIMD}}\text{ (por ciento)}=\frac{S}{8}\times100.
  \label{eq:simd-efficiency}
\end{equation}
La Tabla~\ref{tab:tiempos} presenta los resultados de las 30 repeticiones. Su
número de columnas requiere utilizar el ancho completo de la página.

\begin{table*}[!t]
  \centering
  \caption{Tiempo promedio, desviación estándar y aceleración por tamaño
    ($R=30$).}
  \label{tab:tiempos}
  \small
  \setlength{\tabcolsep}{4pt}
  \begin{tabularx}{\textwidth}{@{}r>{\raggedright\arraybackslash}Xcccc@{}}
    \toprule
    $N$ & Nivel jerárquico y huella &
      Escalar $\overline{T}\pm s$ (ms) &
      AVX2 $\overline{T}\pm s$ (ms) & $S$ & $\eta_{\mathrm{SIMD}}$ \\
    \midrule
    1,000 & L1d (\SI{7.8}{\kibi\byte}) &
      $0.0028\pm0.0001$ & $0.0005\pm0.0000$ & $5.80\times$ & 72.5\% \\
    100,000 & L2 (\SI{781.2}{\kibi\byte}) &
      $0.2943\pm0.0042$ & $0.0525\pm0.0008$ & $5.60\times$ & 70.0\% \\
    1,000,000 & L3 (\SI{7.63}{\mebi\byte}) &
      $2.6637\pm0.0381$ & $0.5433\pm0.0092$ & $4.90\times$ & 61.3\% \\
    20,000,000 & DRAM (\SI{152.6}{\mebi\byte}) &
      $54.5323\pm0.8120$ & $22.3774\pm0.3540$ & $2.44\times$ & 30.5\% \\
    \bottomrule
  \end{tabularx}
\end{table*}

\subsection{Aceleración obtenida}

El límite ideal para AVX2 sobre \texttt{float32} es $8.0\times$, pues cada
registro de 256 bits contiene ocho carriles de 32 bits. La curva de la
Figura~\ref{fig:speedup} muestra dos regímenes. Para
$N\in[10^3,10^5]$ se observa una meseta de $5.80\times$ a $5.60\times$; al
alcanzar $N=10^6$, la aceleración disminuye moderadamente a $4.90\times$ dentro
de la LLC. Cuando la huella rebasa ampliamente la caché, la curva cae y converge
a $2.44\times$ para $N=2\times10^7$, debido al límite de transferencia desde
DRAM.

\begin{figure}[H]
  \centering
  \includegraphics[width=\columnwidth]{docs/speedup_vs_n.png}
  \caption{Aceleración de la implementación AVX2 respecto de la versión escalar.}
  \label{fig:speedup}
\end{figure}

\subsection{Resultados de \texttt{perf stat}}

Los contadores de la Tabla~\ref{tab:perf} se obtuvieron mediante la PMU de
Linux con afinidad al núcleo P0. Se comparan un problema residente en L3
($N=10^6$) y otro dominado por accesos a DRAM ($N=2\times10^7$).

\begin{table*}[!t]
  \centering
  \caption{Contadores microarquitectónicos obtenidos mediante
    \texttt{perf stat} en el núcleo P0 Golden Cove.}
  \label{tab:perf}
  \scriptsize
  \setlength{\tabcolsep}{3pt}
  \begin{tabularx}{\textwidth}{@{}>{\raggedright\arraybackslash}Xrrrr@{}}
    \toprule
    Métrica &
      \multicolumn{1}{c}{$N=10^6$ L3, escalar} &
      \multicolumn{1}{c}{$N=10^6$ L3, AVX2} &
      \multicolumn{1}{c}{$N=2\times10^7$ DRAM, escalar} &
      \multicolumn{1}{c}{$N=2\times10^7$ DRAM, AVX2} \\
    \midrule
    Ciclos & 359,663,277 & 78,283,929 & 7,410,341,145 & 3,356,712,357 \\
    Instrucciones & 916,136,472 & 127,443,561 & 18,313,758,611 & 2,552,451,534 \\
    IPC & 2.55 & 1.63 & 2.47 & 0.76 \\
    Referencias de caché & 10,091,746 & 9,981,784 & 200,364,360 & 199,071,098 \\
    Fallos de caché & 1,918,819 & 915,922 & 140,654,647 & 181,600,111 \\
    Tasa de fallos & 19.01\% & 9.18\% & 70.20\% & 91.22\% \\
    \bottomrule
  \end{tabularx}
\end{table*}

\subsection{Interpretación de resultados}

\subsubsection{Régimen limitado por cómputo}

Para $N\leq10^6$, los datos permanecen en L1d, L2 o L3 y la aceleración se
mantiene entre $5.80\times$ y $4.90\times$. En $N=10^6$, AVX2 retira
127.4 millones de instrucciones frente a 916.1 millones de la versión escalar,
una reducción de $7.19\times$. El empaquetamiento de ocho valores por
instrucción descongestiona la etapa de decodificación y el búfer de reordenamiento.

El IPC escalar resulta mayor que el vectorial (2.55 frente a 1.63), pero esta
métrica cuenta instrucciones, no operaciones útiles por carril. Una instrucción
AVX2 retirada puede realizar ocho operaciones de punto flotante, por lo que un
IPC menor no implica menor rendimiento numérico. El tiempo medido confirma que
la versión vectorial completa el trabajo $4.90\times$ más rápido en L3.

La distancia respecto al límite de $8\times$ se explica por el coste de las
reducciones horizontales, las conversiones escalares, el control de bucles y
los prólogos y epílogos de la ABI. Si se modela la aceleración máxima observada
mediante la ley de Amdahl con una fracción paralelizable $p\approx0.9458$, se
obtiene
\begin{equation}
  S=\frac{1}{(1-p)+p/8}
   =\frac{1}{0.0542+0.9458/8}
   \approx5.80\times.
  \label{eq:amdahl-performance}
\end{equation}
Este modelo atribuye aproximadamente un 5.42\% del tiempo a trabajo no
beneficiado por el ancho vectorial ideal.

\subsubsection{Régimen limitado por memoria}

Para $N=2\times10^7$, los arreglos de entrada y salida ocupan en conjunto
\SI{152.6}{\mebi\byte}, más de quince veces la capacidad de la caché L3 de
\SI{10}{\mebi\byte}. La CPU debe solicitar repetidamente líneas de 64 bytes a
la memoria principal. Aunque las unidades AVX2 consumen 32 bytes por carga, la
tasa de cómputo deja de ser el factor dominante y el rendimiento queda limitado
por el ancho de banda y la latencia de DRAM.

La tasa de fallos medida para AVX2 pasa de 9.18\% en L3 a 91.22\% en el caso de
DRAM, mientras que su IPC cae de 1.63 a 0.76. Las instrucciones vectoriales
continúan reduciendo el número total de instrucciones en aproximadamente
$7.17\times$, pero gran parte de los ciclos se consume esperando datos. Esta
\textit{Memory Wall} limita el speedup temporal a $2.44\times$, demostrando que
el paralelismo aritmético no puede compensar por sí solo una alimentación de
datos insuficiente.

\section{Evidencia de la sesión de GDB}

\subsection{Propósito técnico y microarquitectónico}

La sesión de depuración con GNU GDB auditó en tiempo de ejecución el
cumplimiento de la convención System V AMD64 ABI y contrastó el datapath escalar
SISD con el datapath SIMD AVX2. En particular, se verificó el procesamiento de
un único \texttt{float32} en el carril inferior de \texttt{xmm3}, frente al
procesamiento simultáneo de ocho valores en \texttt{ymm0}. También se comprobó
la alineación a 32 bytes requerida por \texttt{vmovaps} y la igualdad de los
resultados almacenados por ambas implementaciones para el caso $N=16$
(\texttt{tc-06\_in.dat}).

\subsection{Secuencia de comandos y breakpoints}

\subsubsection{Kernel escalar}

La sesión de \texttt{bin/norm\_scalar} se configuró con una repetición para
aislar el flujo de control. El primer breakpoint se ubicó en la entrada de
\texttt{normalize\_array}, sobre \texttt{test edx, edx}, para comprobar los
argumentos \texttt{rdi=in}, \texttt{rsi=out} y \texttt{edx=16}. También se
inspeccionaron $\mu=-5.637640$ en \texttt{xmm0} y
$\sigma=48.620003$ en \texttt{xmm1}.

El segundo breakpoint se colocó inmediatamente después de
\texttt{divss xmm3, xmm1}, dentro de \texttt{.norm\_loop}. Allí se confirmó que
solo el carril inferior de \texttt{xmm3} contenía el resultado de la iteración
actual y que \texttt{eax} avanzaba de uno en uno, equivalente a cuatro bytes
por elemento. Finalmente, se eliminó el breakpoint iterativo y se utilizó
\texttt{finish} antes de inspeccionar la memoria de salida.

\begin{lstlisting}[caption={Comandos principales de la sesión GDB escalar.},label={lst:gdb-scalar}]
(gdb) file bin/norm_scalar
(gdb) set args data/test_suite/tc-06_in.dat data/test_suite/tc-06_sc_out.dat 1
(gdb) break asm/scalar/stats_scalar.asm:167
(gdb) run
(gdb) print $rdi
(gdb) print $rsi
(gdb) print $edx
$edx = 16
(gdb) print $xmm0.v4_float
# carril 0: mean = -5.637640
(gdb) print $xmm1.v4_float
# carril 0: stddev = 48.620003
(gdb) break asm/scalar/stats_scalar.asm:184
(gdb) continue
(gdb) print $eax
$eax = 0
(gdb) print $xmm3.v4_float
# un resultado activo en el carril 0
(gdb) delete 2
(gdb) finish
(gdb) print (void*)out
(gdb) x/8fw out
\end{lstlisting}

\begin{figure}[H]
  \centering
  \IfFileExists{prism-uploads/CAPTURA ESCALAR 1.png}
    {\includegraphics[width=\columnwidth]{prism-uploads/CAPTURA ESCALAR 1.png}}
    {\fbox{\parbox[c][3.0cm][c]{0.92\columnwidth}{\centering
      Captura escalar 1 pendiente: argumentos y valores de
      \texttt{xmm0}/\texttt{xmm1} al entrar en \texttt{normalize\_array}.}}}
  \caption{Recepción de argumentos y parámetros flotantes en el kernel escalar.}
  \label{fig:gdb-escalar-1}
\end{figure}

\begin{figure}[H]
  \centering
  \IfFileExists{prism-uploads/CAPTURA ESCALAR 2.png}
    {\includegraphics[width=\columnwidth]{prism-uploads/CAPTURA ESCALAR 2.png}}
    {\fbox{\parbox[c][3.0cm][c]{0.92\columnwidth}{\centering
      Captura escalar 2 pendiente: \texttt{xmm3} después de
      \texttt{divss}, con un único carril activo y \texttt{eax=0}.}}}
  \caption{Procesamiento SISD de un elemento \texttt{float32}.}
  \label{fig:gdb-escalar-2}
\end{figure}

\begin{figure}[H]
  \centering
  \IfFileExists{prism-uploads/CAPTURA ESCALAR 3.png}
    {\includegraphics[width=\columnwidth]{prism-uploads/CAPTURA ESCALAR 3.png}}
    {\fbox{\parbox[c][3.0cm][c]{0.92\columnwidth}{\centering
      Captura escalar 3 pendiente: volcado \texttt{x/8fw out} después de
      retornar al controlador.}}}
  \caption{Memoria de salida producida por la versión escalar.}
  \label{fig:gdb-escalar-3}
\end{figure}

\subsubsection{Kernel vectorial}

La sesión de \texttt{bin/norm\_vector} utilizó el mismo arreglo y una única
repetición. El breakpoint de entrada, situado sobre \texttt{test edx, edx},
permitió validar los argumentos antes de que \texttt{vbroadcastss} replicara
$\mu$ y $\sigma$ en \texttt{ymm4} y \texttt{ymm5}. El segundo breakpoint se
ubicó después de \texttt{vdivps ymm0, ymm0, ymm5}; en ese punto,
\texttt{ymm0} contenía ocho resultados normalizados y \texttt{eax} avanzaba en
bloques de ocho elementos, equivalentes a 32 bytes.

La expresión \texttt{((unsigned long)out) \% 32} produjo cero, confirmando que
el búfer de salida satisface la alineación exigida por las cargas y
almacenamientos con \texttt{vmovaps}. Tras eliminar el breakpoint del bucle,
\texttt{finish} permitió inspeccionar el mismo intervalo de memoria que en la
versión escalar.

\begin{lstlisting}[caption={Comandos principales de la sesión GDB vectorial.},label={lst:gdb-vector}]
(gdb) file bin/norm_vector
(gdb) set args data/test_suite/tc-06_in.dat data/test_suite/tc-06_vc_out.dat 1
(gdb) break asm/vector/stats_vector.asm:248
(gdb) run
(gdb) print $edx
$edx = 16
(gdb) print $xmm0.v4_float
# mean = -5.637640
(gdb) print $xmm1.v4_float
# stddev = 48.620003
(gdb) break asm/vector/stats_vector.asm:271
(gdb) continue
(gdb) print $eax
$eax = 0
(gdb) print $ymm0.v8_float
# ocho resultados normalizados simultaneos
(gdb) print ((unsigned long)out) % 32
$ = 0
(gdb) delete 2
(gdb) finish
(gdb) x/8fw out
\end{lstlisting}

\begin{figure}[H]
  \centering
  \IfFileExists{prism-uploads/CAPTURA VECTORIAL 1.png}
    {\includegraphics[width=\columnwidth]{prism-uploads/CAPTURA VECTORIAL 1.png}}
    {\fbox{\parbox[c][3.0cm][c]{0.92\columnwidth}{\centering
      Captura vectorial 1 pendiente: argumentos de
      \texttt{normalize\_array} antes de los broadcasts.}}}
  \caption{Argumentos de entrada de la rutina vectorial.}
  \label{fig:gdb-vectorial-1}
\end{figure}

\begin{figure}[H]
  \centering
  \IfFileExists{prism-uploads/CAPTURA VECTORIAL 2.png}
    {\includegraphics[width=\columnwidth]{prism-uploads/CAPTURA VECTORIAL 2.png}}
    {\fbox{\parbox[c][3.0cm][c]{0.92\columnwidth}{\centering
      Captura vectorial 2 pendiente: ocho carriles de \texttt{ymm0} después de
      \texttt{vdivps}.}}}
  \caption{Procesamiento SIMD de ocho elementos \texttt{float32}.}
  \label{fig:gdb-vectorial-2}
\end{figure}

\begin{figure}[H]
  \centering
  \IfFileExists{prism-uploads/CAPTURA VECTORIAL 3.png}
    {\includegraphics[width=\columnwidth]{prism-uploads/CAPTURA VECTORIAL 3.png}}
    {\fbox{\parbox[c][3.0cm][c]{0.92\columnwidth}{\centering
      Captura vectorial 3 pendiente: alineación módulo 32 y volcado
      \texttt{x/8fw out}.}}}
  \caption{Alineación y memoria de salida de la versión AVX2.}
  \label{fig:gdb-vectorial-3}
\end{figure}

\subsection{Comparación de hallazgos en bajo nivel}

La Tabla~\ref{tab:gdb-comparison} conserva el ancho de una columna y resume las
diferencias observadas directamente en GDB.

\begin{table}[H]
  \centering
  \caption{Hallazgos de la inspección escalar y vectorial.}
  \label{tab:gdb-comparison}
  \scriptsize
  \setlength{\tabcolsep}{2.5pt}
  \begin{tabularx}{\columnwidth}{@{}>{\raggedright\arraybackslash}p{0.25\columnwidth}
    >{\raggedright\arraybackslash}X
    >{\raggedright\arraybackslash}X@{}}
    \toprule
    Dimensión & Escalar & AVX2 \\
    \midrule
    Paradigma & SISD, SSE escalar & SIMD empaquetado \\
    Registro & \texttt{xmm3}, carril 0 & \texttt{ymm0}, ocho carriles \\
    Elementos/iteración & 1 & 8 \\
    Paso del índice & $+1$ elemento ($+4$ B) & $+8$ elementos ($+32$ B) \\
    Aritmética & \texttt{subss}, \texttt{divss} & \texttt{vsubps},
      \texttt{vdivps} \\
    Acceso & \texttt{movss}, 32 bits & \texttt{vmovaps}, 256 bits alineados \\
    Parámetros & $\mu$ y $\sigma$ en XMM & Broadcast a \texttt{ymm4} y
      \texttt{ymm5} \\
    Iteraciones ($N=16$) & 16 & 2 \\
    Alineación & 4 bytes suficiente & 32 bytes en esta implementación \\
    Salida inicial & $[0.899751,1.823481,\ldots]$ &
      $[0.899751,1.823481,\ldots]$ \\
    Epílogo & Sin limpieza AVX & \texttt{vzeroupper} \\
    \bottomrule
  \end{tabularx}
\end{table}

La observación confirma una reducción del 87.5\% en el número de iteraciones
para $N=16$. AVX2 ocupa los ocho carriles de un registro YMM y reduce la
sobrecarga de control, mientras que la versión escalar utiliza únicamente un
valor de 32 bits por instrucción. La alineación a 32 bytes evita accesos
desalineados con \texttt{vmovaps}; \texttt{vzeroupper}, por su parte, prepara el
retorno a código que pueda utilizar SSE heredado.

\subsection{Convergencia numérica observada}

La inspección \texttt{x/8fw out} produjo en ambos ejecutables el mismo bloque
inicial:
\begin{equation}
\begin{aligned}
  \mathrm{out}[0\ldots7]=[&0.899751365,\;1.823480730,\\
  &-0.0602805801,\;-0.851318777,\\
  &-1.114069460,\;-1.663971660,\\
  &0.220615491,\;0.604119360].
\end{aligned}
\label{eq:gdb-output-values}
\end{equation}
Para estos ocho elementos, la diferencia máxima observada fue
\begin{equation}
  \max_{0\leq i<8}
  \left|\mathrm{out}_{\mathrm{sc}}[i]
  -\mathrm{out}_{\mathrm{vec}}[i]\right|=0.0,
  \label{eq:gdb-max-difference}
\end{equation}
valor inferior a la epsilon de máquina de \texttt{float32},
$\varepsilon\approx1.19\times10^{-7}$. La evidencia confirma la coincidencia
exacta del bloque inspeccionado, además de la equivalencia numérica global
establecida por la suite de pruebas.

\section{Conclusiones, limitaciones y trabajo futuro}

\subsection{Conclusiones}

La vectorización mediante AVX2 de 256 bits demostró la eficacia del paralelismo
a nivel de datos para el cálculo de estadísticos y la normalización de arreglos
\texttt{float32}. La implementación alcanzó un speedup máximo de
$5.80\times$, equivalente al 72.5\% del límite ideal de $8.0\times$, y redujo
en $7.19\times$ el número de instrucciones retiradas en el caso residente en
L3. Esta mejora se obtuvo sin comprometer la exactitud: el error relativo máximo
de la suite fue $1.19\times10^{-6}$, muy inferior a la tolerancia establecida de
$1.0\times10^{-4}$.

Los resultados también evidencian la transición entre dos regímenes de
ejecución. Mientras los datos permanecen en la jerarquía de cachés, para
$N\leq10^6$, la aplicación se encuentra principalmente limitada por cómputo y
alcanza aceleraciones entre $5.80\times$ y $4.90\times$. Al utilizar
$N=2\times10^7$, cuya huella combinada es de \SI{152.6}{\mebi\byte}, el
problema pasa a estar limitado por memoria. La tasa de fallos de caché de la
versión AVX2 alcanza 91.22\% y el speedup desciende a $2.44\times$, confirmando
el efecto de la \textit{Memory Wall}.

La ley de Amdahl explica la distancia restante respecto al factor ideal de
ocho. Las reducciones horizontales con \texttt{vextractf128} y
\texttt{vhaddps}, las conversiones escalares, el control de los bucles y el
procesamiento del remanente cuando $N\bmod8\neq0$ constituyen una fracción que
no obtiene el beneficio completo del ancho SIMD. Este coste fijo tiene un peso
relativo mayor en arreglos pequeños.

Finalmente, la implementación mostró robustez en el nivel de arquitectura y de
interfaz binaria. El cumplimiento de System V AMD64 ABI permitió integrar los
núcleos NASM con el controlador en C; la asignación alineada a 32 bytes evitó
accesos inválidos con \texttt{vmovaps}; y \texttt{vzeroupper} dejó limpios los
carriles superiores antes de retornar a código potencialmente basado en SSE.
La inspección con GDB confirmó tanto la alineación como la equivalencia de las
salidas escalares y vectoriales.

\subsection{Limitaciones observadas}

La limitación dominante para arreglos que exceden la LLC es el ancho de banda
de DRAM. En este régimen, la capacidad de las unidades vectoriales para consumir
y procesar datos supera la tasa con la que la memoria principal puede entregar
líneas de caché. En consecuencia, las unidades SIMD permanecen parcialmente
inactivas mientras el pipeline espera datos, por lo que aumentar únicamente el
ancho vectorial no produce una aceleración proporcional.

La implementación también es estrictamente monohilo y utiliza un solo núcleo
físico. Por ello, no aprovecha el paralelismo a nivel de hilos ni los restantes
P-cores y E-cores del procesador. Los resultados caracterizan el beneficio de
SIMD dentro de un núcleo, pero no representan el rendimiento máximo agregado de
la plataforma completa.

Adicionalmente, las sumas se realizan en precisión simple y su orden cambia
entre la versión escalar y la vectorial. Aunque las diferencias permanecieron
dentro de la tolerancia definida, aplicaciones con requisitos numéricos más
estrictos podrían necesitar acumuladores de doble precisión o algoritmos de
suma compensada.

\subsection{Trabajo futuro}

Una primera extensión consiste en emplear FMA3 durante la acumulación de la
varianza. Después de calcular $d_i=x_i-\mu$, una instrucción como
\texttt{vfmadd231ps} puede evaluar y acumular $d_i^2$ en una única operación
fusionada. Esto reduce el número de instrucciones y evita el redondeo intermedio
separado entre la multiplicación y la suma.

También se propone desenrollar los bucles vectoriales por factores de dos o
cuatro y mantener varios acumuladores YMM independientes. Esta organización
reduce las dependencias RAW de las cadenas de suma y ofrece al planificador
fuera de orden más operaciones independientes para utilizar los puertos de
ejecución disponibles.

Para arreglos mucho mayores que la LLC, podría evaluarse el uso de
almacenamientos no temporales mediante \texttt{vmovntps}. Si la salida no será
reutilizada inmediatamente, estas escrituras pueden disminuir el tráfico de
\textit{write-allocate} y reducir la contaminación de las cachés. Su beneficio
debe medirse experimentalmente, pues depende del tamaño, la alineación y el
patrón de reutilización.

Por último, una implementación híbrida SIMD+OpenMP permitiría dividir el arreglo
entre varios hilos y combinar el paralelismo AVX2 dentro de cada núcleo con
paralelismo a nivel de hilos. Esta alternativa deberá considerar la topología
híbrida del procesador, la afinidad de los hilos y la saturación compartida del
ancho de banda de memoria, que probablemente se convierta nuevamente en el
límite para tamaños grandes.

\clearpage
\appendices
\onecolumn
\section{Diagramas de arquitectura}

La Figura~\ref{fig:arquitectura-software-hardware} resume el flujo jerárquico
del sistema, desde la orquestación en C hasta la ejecución de los kernels
escalares y vectoriales en los registros del procesador.

\begin{figure}[H]
  \centering
  \resizebox{0.92\textwidth}{!}{%
  \begin{tikzpicture}[
    font=\sffamily,
    flow/.style={-{Latex[length=3mm,width=2mm]},thick,draw=black!65},
    mainnode/.style={
      rectangle,rounded corners=2pt,draw=blue!65!black,fill=blue!8,
      thick,align=center,minimum height=9mm,text width=47mm
    },
    cnode/.style={
      rectangle,rounded corners=2pt,draw=blue!55!black,fill=blue!4,
      align=center,minimum height=10mm,text width=42mm
    },
    inode/.style={
      rectangle,rounded corners=2pt,draw=teal!65!black,fill=teal!7,
      align=center,minimum height=12mm,text width=50mm
    },
    abinode/.style={
      rectangle,rounded corners=2pt,draw=orange!75!black,fill=orange!9,
      align=center,minimum height=14mm,text width=74mm
    },
    scalar/.style={
      rectangle,rounded corners=2pt,draw=violet!70!black,fill=violet!8,
      align=center,minimum height=13mm,text width=73mm
    },
    vector/.style={
      rectangle,rounded corners=2pt,draw=green!50!black,fill=green!9,
      align=center,minimum height=13mm,text width=73mm
    },
    group/.style={
      rectangle,rounded corners=4pt,thick,inner sep=6mm
    },
    grouptitle/.style={font=\bfseries\sffamily,align=center}
  ]

  % Capa de orquestación en C.
  \node[mainnode] (cmain) {\textbf{\texttt{main()}}\\Orquestador del benchmark};
  \node[cnode,below left=10mm and 15mm of cmain] (calloc)
    {\texttt{aligned\_alloc}\\\texttt{(32, size)}\\Reserva alineada en el heap};
  \node[cnode,below=10mm of cmain] (cbench)
    {\texttt{clock\_gettime}\\Telemetría temporal};
  \node[cnode,below right=10mm and 15mm of cmain] (cmath)
    {\texttt{sqrtf(var)}\\$\sigma=\sqrt{\mathrm{var}}$};
  \draw[flow] (cmain) -- (calloc);
  \draw[flow] (cmain) -- (cbench);
  \draw[flow] (cmain) -- (cmath);

  % Interfaz pública.
  \node[inode,below=27mm of calloc] (hsum)
    {\texttt{sum\_array(const float *arr, int n)}};
  \node[inode,below=27mm of cbench] (hstats)
    {\texttt{compute\_stats(...)}\\Media, varianza, mínimo y máximo};
  \node[inode,below=27mm of cmath] (hnorm)
    {\texttt{normalize\_array(...)}\\Normalización \textit{z-score}};

  % Contrato System V AMD64 ABI.
  \node[abinode,below left=27mm and 5mm of hstats] (abigpr)
    {\textbf{Punteros y enteros en GPR}\\
     1: RDI $\mid$ 2: RSI $\mid$ 3: RDX $\mid$ 4: RCX $\mid$ 5: R8 $\mid$ 6: R9};
  \node[abinode,below right=27mm and 5mm of hstats] (abixmm)
    {\textbf{Flotantes escalares y retorno}\\
     1.er float: XMM0 $\mid$ 2.º float: XMM1 $\mid$ retorno: XMM0};
  \node[abinode,below=9mm of abigpr] (abipreserve)
    {\textbf{Registros no volátiles}\\
     RBX, RBP y R12--R15 preservados con \texttt{push}/\texttt{pop};
     pila alineada a 16 bytes};
  \node[abinode,below=9mm of abixmm] (abiclean)
    {\textbf{Higiene del estado AVX}\\
     \texttt{vzeroupper} antes de \texttt{ret} para evitar transiciones
     AVX--SSE costosas};

  % Backends escalares y vectoriales.
  \node[scalar,below=31mm of abipreserve] (ssum)
    {\texttt{sum\_array}\\SISD secuencial: \texttt{addss}, un float/iteración};
  \node[scalar,below=7mm of ssum] (sstats)
    {\texttt{compute\_stats}\\Dos pasadas: \texttt{minss}, \texttt{maxss},
     \texttt{divss}};
  \node[scalar,below=7mm of sstats] (snorm)
    {\texttt{normalize\_array}\\Traslación y escalado elemento a elemento};

  \node[vector,below=31mm of abiclean] (vsum)
    {\texttt{sum\_array}\\\texttt{vmovaps}/\texttt{vaddps}, ocho floats/iteración\\
     Reducción horizontal, \textit{tail loop} y \texttt{vzeroupper}};
  \node[vector,below=7mm of vsum] (vstats)
    {\texttt{compute\_stats}\\Dos pasadas AVX2 con \texttt{vminps} y
     \texttt{vmaxps}\\Árbol de reducción, \textit{tail loop} y \texttt{vzeroupper}};
  \node[vector,below=7mm of vstats] (vnorm)
    {\texttt{normalize\_array}\\Broadcast de $\mu$ y $\sigma$;
     \texttt{vsubps}/\texttt{vdivps}\\\textit{Tail loop} y
     \texttt{vzeroupper}};

  % Contenedores de las capas, dibujados detrás de los nodos.
  \begin{scope}[on background layer]
    \node[group,draw=blue!65!black,fill=blue!2,
      fit=(cmain)(calloc)(cbench)(cmath),
      label={[grouptitle,blue!65!black]above:
        Capa de orquestación en C: \texttt{src/driver.c}}] (cgroup) {};

    \node[group,draw=teal!65!black,fill=teal!2,
      fit=(hsum)(hstats)(hnorm),
      label={[grouptitle,teal!65!black]above:
        Interfaz y contrato de software: \texttt{include/stats.h}}] (igroup) {};

    \node[group,draw=orange!75!black,fill=orange!3,
      fit=(abigpr)(abixmm)(abipreserve)(abiclean),
      label={[grouptitle,orange!75!black]above:
        Protocolo de hardware: System V AMD64 ABI}] (abigroup) {};

    \node[group,draw=violet!70!black,fill=violet!2,
      fit=(ssum)(sstats)(snorm),
      label={[grouptitle,violet!70!black]above:
        Backend escalar: \texttt{asm/scalar/stats\_scalar.asm}}] (sgroup) {};

    \node[group,draw=green!50!black,fill=green!3,
      fit=(vsum)(vstats)(vnorm),
      label={[grouptitle,green!50!black]above:
        Backend AVX2: \texttt{asm/vector/stats\_vector.asm}}] (vgroup) {};

    \node[group,draw=black!60,fill=black!1,inner sep=10mm,
      fit=(sgroup)(vgroup),
      label={[grouptitle]above:Kernels en ensamblador NASM x86-64}] (kgroup) {};
  \end{scope}

  % Flujo jerárquico entre capas.
  \draw[flow] (cgroup.south) -- (igroup.north);
  \draw[flow] (igroup.south) -- (abigroup.north);
  \draw[flow] (abigroup.south) -- ++(0,-7mm) -| (sgroup.north);
  \draw[flow] (abigroup.south) -- ++(0,-7mm) -| (vgroup.north);

  \end{tikzpicture}%
  }
  \caption{Flujo jerárquico desde la capa de orquestación en C hasta los
    backends escalar y AVX2 ejecutados por el hardware.}
  \label{fig:arquitectura-software-hardware}
\end{figure}

\clearpage
\subsection{Bucle escalar SISD paso a paso}

El kernel escalar recorre el arreglo de manera estrictamente secuencial. Cada
iteración carga un único valor \texttt{float32}, lo acumula en el carril
inferior de \texttt{xmm0} y avanza el índice una posición. Al no existir
particionamiento vectorial, la rutina no necesita calcular un límite por bloques
ni realizar una reducción horizontal al finalizar.

\begin{figure}[H]
  \centering
  \resizebox{0.82\textwidth}{!}{%
  \begin{tikzpicture}[
    font=\sffamily,
    node distance=10mm,
    flow/.style={-{Latex[length=3mm,width=2mm]},thick,draw=black!70},
    terminal/.style={
      ellipse,draw=blue!65!black,fill=blue!8,thick,
      align=center,minimum height=12mm,text width=62mm
    },
    process/.style={
      rectangle,rounded corners=2pt,draw=teal!65!black,fill=teal!7,
      thick,align=center,minimum height=14mm,text width=74mm
    },
    decision/.style={
      diamond,aspect=2.1,draw=orange!75!black,fill=orange!10,
      thick,align=center,inner sep=1.5mm,text width=36mm
    },
    branch/.style={font=\bfseries\small,fill=white,inner sep=1pt}
  ]

  \node[terminal] (start)
    {\textbf{Inicio de \texttt{sum\_array} escalar (SISD)}\\
     \texttt{rdi = arr}, \texttt{esi = n}};

  \node[process,below=of start] (init)
    {\textbf{Inicializar registros}\\
     \texttt{eax = 0} \quad ($i=0$)\\
     \texttt{xorps xmm0, xmm0} \quad ($\mathrm{sum}=0.0f$)};

  \node[decision,below=12mm of init] (cond)
    {$i<n$?\\\texttt{cmp eax, esi}\\\texttt{jl}};

  \node[process,below=13mm of cond] (load)
    {\textbf{Cargar un valor \texttt{float32}}\\
     \texttt{movss xmm1, [rdi + rax*4]}};

  \node[process,below=of load] (arith)
    {\textbf{Acumular la suma escalar}\\
     \texttt{addss xmm0, xmm1}};

  \node[process,below=of arith] (inc)
    {\textbf{Avanzar el índice}\\
     \texttt{inc eax} \quad ($i\mathrel{+}=1$)};

  \node[terminal,right=38mm of cond] (ret)
    {\textbf{Retornar}\\
     \texttt{ret}\\Resultado escalar en \texttt{xmm0}};

  \draw[flow] (start) -- (init);
  \draw[flow] (init) -- (cond);
  \draw[flow] (cond) -- node[branch,left]{Sí} (load);
  \draw[flow] (load) -- (arith);
  \draw[flow] (arith) -- (inc);
  \draw[flow] (cond) -- node[branch,above]{No} (ret);
  \draw[flow] (inc.east) -- ++(28mm,0) |- (cond.east);

  \end{tikzpicture}%
  }
  \caption{Flujo de ejecución del bucle escalar SISD de
    \texttt{sum\_array}.}
  \label{fig:flujo-sum-array-escalar}
\end{figure}

\clearpage
\subsection{Bucle vectorial AVX2 con remanente}

El kernel vectorial procesa ocho valores \texttt{float32} por iteración. Antes
de entrar al bucle calcula $N_{\mathrm{vec}}=n\mathbin{\&}(\mathord{\sim}7)$,
reduce horizontalmente los acumuladores al concluir los bloques completos y
procesa los $n\bmod8$ elementos restantes mediante un bucle escalar de cierre.

\begin{figure}[H]
  \centering
  \resizebox{!}{0.72\textheight}{%
  \begin{tikzpicture}[
    font=\sffamily,
    node distance=7mm,
    flow/.style={-{Latex[length=3mm,width=2mm]},thick,draw=black!70},
    terminal/.style={
      ellipse,draw=blue!65!black,fill=blue!8,thick,
      align=center,minimum height=11mm,text width=76mm
    },
    process/.style={
      rectangle,rounded corners=2pt,draw=green!50!black,fill=green!8,
      thick,align=center,minimum height=12mm,text width=82mm
    },
    reduction/.style={
      rectangle,rounded corners=2pt,draw=violet!70!black,fill=violet!8,
      thick,align=left,minimum height=25mm,text width=100mm
    },
    cleanup/.style={
      rectangle,rounded corners=2pt,draw=orange!75!black,fill=orange!9,
      thick,align=center,minimum height=12mm,text width=82mm
    },
    decision/.style={
      diamond,aspect=2.25,draw=orange!75!black,fill=orange!10,
      thick,align=center,inner sep=1.3mm,text width=42mm
    },
    branch/.style={font=\bfseries\small,fill=white,inner sep=1pt}
  ]

  \node[terminal] (vstart)
    {\textbf{Inicio de \texttt{sum\_array} vectorial (AVX2)}\\
     \texttt{rdi = arr} alineado a 32 B, \texttt{esi = n}};

  \node[process,below=of vstart] (vinit)
    {\textbf{Inicializar índice y acumulador}\\
     \texttt{eax = 0} \quad ($i=0$)\\
     \texttt{vxorps ymm0, ymm0, ymm0} \quad (ocho sumas en cero)};

  \node[process,below=of vinit] (vlimit)
    {\textbf{Calcular el límite vectorial}\\
     \texttt{ecx = esi \& \textasciitilde 7}\\
     $N_{\mathrm{vec}}$ es múltiplo de ocho};

  \node[decision,below=10mm of vlimit] (vcheck)
    {$\texttt{ecx}\geq8$?\\¿Existe un bloque completo?};

  \node[decision,below=11mm of vcheck] (vcond)
    {$i<\texttt{ecx}$?\\\texttt{cmp eax, ecx} / \texttt{jl}};

  \node[process,below=11mm of vcond] (vload)
    {\textbf{Carga alineada de 256 bits}\\
     \texttt{vmovaps ymm1, [rdi + rax*4]}};

  \node[process,below=of vload] (vadd)
    {\textbf{Suma empaquetada en ocho carriles}\\
     \texttt{vaddps ymm0, ymm0, ymm1}};

  \node[process,below=of vadd] (vinc)
    {\textbf{Avanzar el índice vectorial}\\
     \texttt{add eax, 8} \quad ($i\mathrel{+}=8$)};

  \node[reduction,right=34mm of vcond] (vreduce)
    {\textbf{Reducción horizontal en árbol: $8\rightarrow1$}\\[1mm]
     1. \texttt{vextractf128 xmm2, ymm0, 1}\\
     2. \texttt{vaddps xmm0, xmm0, xmm2}\\
     3. \texttt{vhaddps xmm0, xmm0, xmm0}\\
     4. \texttt{vhaddps xmm0, xmm0, xmm0}\\
     Resultado total en \texttt{xmm0[0]}};

  \node[decision,below=30mm of vinc] (tailcond)
    {$i<\texttt{esi}$?\\¿Quedan $n\bmod8$ elementos?};

  \node[process,below=11mm of tailcond] (tailload)
    {\textbf{Carga escalar segura}\\
     \texttt{vmovss xmm1, [rdi + rax*4]}};

  \node[process,below=of tailload] (tailadd)
    {\textbf{Acumular sobre el resultado reducido}\\
     \texttt{vaddss xmm0, xmm0, xmm1}};

  \node[process,below=of tailadd] (tailinc)
    {\textbf{Avanzar el índice escalar}\\
     \texttt{inc eax} \quad ($i\mathrel{+}=1$)};

  \node[cleanup,right=35mm of tailcond] (vzero)
    {\textbf{Limpiar los carriles superiores}\\
     \texttt{vzeroupper}\\Evita transiciones AVX--SSE costosas};

  \node[terminal,below=of vzero] (vret)
    {\textbf{Retornar}\\
     \texttt{ret}; resultado escalar en \texttt{xmm0}};

  % Flujo principal y ramas.
  \draw[flow] (vstart) -- (vinit);
  \draw[flow] (vinit) -- (vlimit);
  \draw[flow] (vlimit) -- (vcheck);
  \draw[flow] (vcheck) -- node[branch,left]{Sí} (vcond);
  \draw[flow] (vcond) -- node[branch,left]{Sí} (vload);
  \draw[flow] (vload) -- (vadd);
  \draw[flow] (vadd) -- (vinc);
  \draw[flow] (vinc.west) -- ++(-24mm,0) |- (vcond.west);
  \draw[flow] (vcond) -- node[branch,above]{No} (vreduce);
  \draw[flow] (vreduce.south) |- (tailcond.east);
  \draw[flow] (vcheck.west) -- ++(-30mm,0) node[branch,above]{No: $n<8$}
    |- (tailcond.west);
  \draw[flow] (tailcond) -- node[branch,left]{Sí} (tailload);
  \draw[flow] (tailload) -- (tailadd);
  \draw[flow] (tailadd) -- (tailinc);
  \draw[flow] (tailinc.west) -- ++(-24mm,0) |- (tailcond.west);
  \draw[flow] (tailcond) -- node[branch,above]{No} (vzero);
  \draw[flow] (vzero) -- (vret);

  \end{tikzpicture}%
  }
  \caption{Flujo del bucle AVX2, la reducción horizontal y el procesamiento
    escalar del remanente.}
  \label{fig:flujo-sum-array-avx2}
\end{figure}

\clearpage
\subsection{Normalización vectorial con guarda de desviación nula}

La rutina \texttt{normalize\_array} calcula
\begin{equation}
  \mathrm{out}[i]=\frac{\mathrm{in}[i]-\mu}{\sigma}.
  \label{eq:appendix-vector-normalization}
\end{equation}
Los parámetros escalares $\mu$ y $\sigma$ se replican en los ocho carriles de
registros YMM mediante \texttt{vbroadcastss}. Antes de efectuar cualquier
división, la rutina comprueba si $\sigma=0.0$; en ese caso ejecuta una copia
directa para evitar valores NaN o infinitos. La Figura~\ref{fig:flujo-normalize-avx2}
presenta ambas rutas y sus bucles escalares de remanente.

\begin{figure}[H]
  \centering
  \resizebox{!}{0.76\textheight}{%
  \begin{tikzpicture}[
    font=\sffamily,
    node distance=7mm,
    flow/.style={-{Latex[length=2.8mm,width=1.9mm]},thick,draw=black!70},
    terminal/.style={
      ellipse,draw=blue!65!black,fill=blue!8,thick,
      align=center,minimum height=10mm,text width=70mm
    },
    guard/.style={
      rectangle,rounded corners=2pt,draw=orange!75!black,fill=orange!9,
      thick,align=center,minimum height=12mm,text width=76mm
    },
    copy/.style={
      rectangle,rounded corners=2pt,draw=violet!70!black,fill=violet!8,
      thick,align=center,minimum height=12mm,text width=76mm
    },
    norm/.style={
      rectangle,rounded corners=2pt,draw=green!50!black,fill=green!8,
      thick,align=center,minimum height=12mm,text width=82mm
    },
    decision/.style={
      diamond,aspect=2.15,draw=orange!75!black,fill=orange!10,
      thick,align=center,inner sep=1.2mm,text width=39mm
    },
    branch/.style={font=\bfseries\small,fill=white,inner sep=1pt}
  ]

  % Entrada y guardas comunes.
  \node[terminal] (nstart)
    {\textbf{Inicio de \texttt{normalize\_array} vectorial}\\
     \texttt{rdi=in}, \texttt{rsi=out}, \texttt{edx=n}\\
     \texttt{xmm0=mean}, \texttt{xmm1=stddev}};

  \node[decision,below=9mm of nstart] (checkn)
    {$n>0$?\\\texttt{test edx, edx} / \texttt{jg}};

  \node[guard,below=10mm of checkn] (zerosigma)
    {\textbf{Comparar la desviación con cero}\\
     \texttt{vxorps xmm2, xmm2, xmm2}\\
     \texttt{vucomiss xmm1, xmm2}};

  \node[decision,below=10mm of zerosigma] (checkzero)
    {$\sigma=0.0f$?\\\texttt{je .norm\_copy\_loop}};

  % Ruta izquierda: copia directa.
  \node[copy,below left=15mm and 31mm of checkzero] (copyprep)
    {\textbf{Caso singular: copia sin modificar}\\
     \texttt{ecx = edx \& \textasciitilde 7}\\
     \texttt{eax = 0}};

  \node[decision,below=10mm of copyprep] (copyveccond)
    {$i<\texttt{ecx}$?\\¿Queda un bloque de ocho?};

  \node[copy,below=10mm of copyveccond] (copyvecbody)
    {\textbf{Copiar ocho valores}\\
     \texttt{vmovaps ymm0, [rdi + rax*4]}\\
     \texttt{vmovaps [rsi + rax*4], ymm0}\\
     \texttt{add eax, 8}};

  \node[decision,below=15mm of copyvecbody] (copytailcond)
    {$i<\texttt{edx}$?\\¿Queda un remanente?};

  \node[copy,below=10mm of copytailcond] (copytailbody)
    {\textbf{Copiar un valor remanente}\\
     \texttt{vmovss xmm0, [rdi + rax*4]}\\
     \texttt{vmovss [rsi + rax*4], xmm0}\\
     \texttt{inc eax}};

  \node[guard,below=12mm of copytailbody] (vzerocopy)
    {\texttt{vzeroupper}};

  \node[terminal,below=of vzerocopy] (retdirect)
    {\textbf{Retornar directamente}\\\texttt{ret}};

  % Ruta derecha: normalización AVX2.
  \node[norm,below right=15mm and 31mm of checkzero] (broadcast)
    {\textbf{Replicar parámetros en 256 bits}\\
     \texttt{vbroadcastss ymm4, xmm0} \quad (ocho medias)\\
     \texttt{vbroadcastss ymm5, xmm1} \quad (ocho desviaciones)};

  \node[norm,below=of broadcast] (normprep)
    {\textbf{Preparar el límite vectorial}\\
     \texttt{ecx = edx \& \textasciitilde 7}\\
     \texttt{eax = 0}};

  \node[decision,below=10mm of normprep] (normveccond)
    {$i<\texttt{ecx}$?\\¿Queda un bloque de ocho?};

  \node[norm,below=10mm of normveccond] (normvecstep)
    {\textbf{Normalizar ocho valores en paralelo}\\
     \texttt{vmovaps ymm0, [rdi + rax*4]}\\
     \texttt{vsubps ymm0, ymm0, ymm4}\\
     \texttt{vdivps ymm0, ymm0, ymm5}\\
     \texttt{vmovaps [rsi + rax*4], ymm0}};

  \node[norm,below=of normvecstep] (normvecinc)
    {\textbf{Avanzar el índice vectorial}\\
     \texttt{add eax, 8}};

  \node[decision,below=10mm of normvecinc] (normtailcond)
    {$i<\texttt{edx}$?\\¿Quedan $n\bmod8$ valores?};

  \node[norm,below=10mm of normtailcond] (normtailstep)
    {\textbf{Normalizar un valor remanente}\\
     \texttt{vmovss xmm2, [rdi + rax*4]}\\
     \texttt{vsubss xmm2, xmm2, xmm4}\\
     \texttt{vdivss xmm2, xmm2, xmm5}\\
     \texttt{vmovss [rsi + rax*4], xmm2}};

  \node[norm,below=of normtailstep] (normtailinc)
    {\textbf{Avanzar el índice escalar}\\
     \texttt{inc eax}};

  \node[guard,below=12mm of normtailinc] (normvzero)
    {\textbf{Limpiar los carriles superiores}\\
     \texttt{vzeroupper}};

  \node[terminal,below=of normvzero] (retfinal)
    {\textbf{Retornar}\\\texttt{ret}};

  % Flujo común.
  \draw[flow] (nstart) -- (checkn);
  \draw[flow] (checkn) -- node[branch,left]{Sí} (zerosigma);
  \draw[flow] (zerosigma) -- (checkzero);
  \draw[flow] (checkn.west) -- ++(-38mm,0) node[branch,above]{No}
    |- (retdirect.west);

  % Rama de copia.
  \draw[flow] (checkzero) -- node[branch,above left]{Sí: $\sigma=0$} (copyprep);
  \draw[flow] (copyprep) -- (copyveccond);
  \draw[flow] (copyveccond) -- node[branch,left]{Sí} (copyvecbody);
  \draw[flow] (copyvecbody.west) -- ++(-22mm,0) |- (copyveccond.west);
  \draw[flow] (copyveccond.east) -- ++(13mm,0) node[branch,above]{No}
    |- (copytailcond.east);
  \draw[flow] (copytailcond) -- node[branch,left]{Sí} (copytailbody);
  \draw[flow] (copytailbody.west) -- ++(-22mm,0) |- (copytailcond.west);
  \draw[flow] (copytailcond.east) -- ++(13mm,0) node[branch,above]{No}
    |- (vzerocopy.east);
  \draw[flow] (vzerocopy) -- (retdirect);

  % Rama de normalización.
  \draw[flow] (checkzero) -- node[branch,above right]{No: $\sigma>0$} (broadcast);
  \draw[flow] (broadcast) -- (normprep);
  \draw[flow] (normprep) -- (normveccond);
  \draw[flow] (normveccond) -- node[branch,left]{Sí} (normvecstep);
  \draw[flow] (normvecstep) -- (normvecinc);
  \draw[flow] (normvecinc.east) -- ++(22mm,0) |- (normveccond.east);
  \draw[flow] (normveccond.west) -- ++(-13mm,0) node[branch,above]{No}
    |- (normtailcond.west);
  \draw[flow] (normtailcond) -- node[branch,left]{Sí} (normtailstep);
  \draw[flow] (normtailstep) -- (normtailinc);
  \draw[flow] (normtailinc.east) -- ++(22mm,0) |- (normtailcond.east);
  \draw[flow] (normtailcond.west) -- ++(-13mm,0) node[branch,above]{No}
    |- (normvzero.west);
  \draw[flow] (normvzero) -- (retfinal);

  \end{tikzpicture}%
  }
  \caption{Flujo de \texttt{normalize\_array}: guarda para $\sigma=0$, copia
    directa y normalización vectorial con bucles de remanente.}
  \label{fig:flujo-normalize-avx2}
\end{figure}

\clearpage
\subsection{Flujo y validación de memoria}

\begin{figure}[H]
  \centering
  \resizebox{0.94\textwidth}{!}{%
  \begin{tikzpicture}[
    font=\sffamily,
    node distance=9mm,
    flow/.style={-{Latex[length=3mm,width=2mm]},thick,draw=black!70},
    terminal/.style={
      ellipse,draw=blue!65!black,fill=blue!8,thick,
      align=center,minimum height=11mm,text width=76mm
    },
    memory/.style={
      rectangle,rounded corners=2pt,draw=blue!60!black,fill=blue!6,
      thick,align=center,minimum height=13mm,text width=92mm
    },
    fault/.style={
      rectangle,rounded corners=2pt,draw=red!70!black,fill=red!8,
      thick,align=center,minimum height=14mm,text width=82mm
    },
    scalar/.style={
      rectangle,rounded corners=2pt,draw=violet!70!black,fill=violet!8,
      thick,align=left,minimum height=22mm,text width=82mm
    },
    vector/.style={
      rectangle,rounded corners=2pt,draw=green!50!black,fill=green!8,
      thick,align=left,minimum height=22mm,text width=82mm
    },
    decision/.style={
      diamond,aspect=2.2,draw=orange!75!black,fill=orange!10,
      thick,align=center,inner sep=1.3mm,text width=42mm
    },
    branch/.style={font=\bfseries\small,fill=white,inner sep=1pt}
  ]

  \node[terminal] (pstart)
    {\textbf{Puntero de memoria \texttt{arr}}\\
     Retornado por \texttt{aligned\_alloc(32, size)}};

  \node[decision,below=10mm of pstart] (pcheck)
    {¿Dirección alineada a 32 bytes?\\
     $\texttt{arr}\mathbin{\&}\texttt{0x1F}=0$};

  \node[fault,left=38mm of pcheck] (gpfault)
    {\textbf{Fallo general de protección \#GP(0)}\\
     La CPU aborta \texttt{vmovaps}\\
     El sistema entrega la señal SIGSEGV};

  \node[terminal,below=of gpfault] (abort)
    {\textbf{Terminación abrupta del programa}};

  \node[memory,below=12mm of pcheck] (valid)
    {\textbf{Dirección válida para \texttt{vmovaps}}\\
     Terminaciones posibles: \texttt{0x...00}, \texttt{0x...20},
     \texttt{0x...40} o \texttt{0x...60}};

  \node[memory,below=of valid] (cacheline)
    {\textbf{Mapeo en una línea L1d de 64 bytes}\\
     Offset \texttt{0x00}: $f_0\ldots f_7$ (32 B)\\
     Offset \texttt{0x20}: $f_8\ldots f_{15}$ (32 B)\\
     Sin división de una carga AVX2 entre dos líneas de caché};

  \node[decision,below=11mm of cacheline] (kernelbranch)
    {Selección del kernel\\de ejecución};

  \node[scalar,below left=15mm and 24mm of kernelbranch] (strideScalar)
    {\textbf{Kernel escalar SISD: avance de 4 B}\\[1mm]
     1. \texttt{movss xmm1, [rdi + rax*4]}\\
     2. \texttt{addss xmm0, xmm1}\\
     3. \texttt{inc eax}\\
     16 iteraciones por línea de caché};

  \node[vector,below right=15mm and 24mm of kernelbranch] (strideVector)
    {\textbf{Kernel vectorial AVX2: avance de 32 B}\\[1mm]
     1. \texttt{vmovaps ymm1, [rdi + rax*4]}\\
     2. \texttt{vaddps ymm0, ymm0, ymm1}\\
     3. \texttt{add eax, 8}\\
     2 iteraciones por línea de caché};

  \node[terminal,below=24mm of kernelbranch] (pend)
    {\textbf{Fin del recorrido del arreglo}};

  \draw[flow] (pstart) -- (pcheck);
  \draw[flow] (pcheck) -- node[branch,above]{No: desalineado} (gpfault);
  \draw[flow] (gpfault) -- (abort);
  \draw[flow] (pcheck) -- node[branch,left]{Sí} (valid);
  \draw[flow] (valid) -- (cacheline);
  \draw[flow] (cacheline) -- (kernelbranch);
  \draw[flow] (kernelbranch) -- node[branch,above left]{Escalar} (strideScalar);
  \draw[flow] (kernelbranch) -- node[branch,above right]{AVX2} (strideVector);
  \draw[flow] (strideScalar.south) |- (pend.west);
  \draw[flow] (strideVector.south) |- (pend.east);

  \end{tikzpicture}%
  }
  \caption{Validación de alineación, mapeo en líneas de caché y patrones de
    recorrido escalar y AVX2.}
  \label{fig:flujo-validacion-memoria}
\end{figure}

\clearpage
\subsection{Diagrama 3: Tablas exhaustivas de asignación de registros}

\begingroup
\scriptsize
\setlength{\tabcolsep}{3pt}
\renewcommand{\arraystretch}{1.18}

\begin{longtable}{@{}>{\raggedright\arraybackslash}p{0.16\textwidth}
  >{\raggedright\arraybackslash}p{0.23\textwidth}
  >{\raggedright\arraybackslash}p{0.20\textwidth}
  >{\raggedright\arraybackslash}p{0.35\textwidth}@{}}
\caption{Convención de registros en System V AMD64 ABI.}
\label{tab:abi-register-convention}\\
\toprule
Clase ABI & Registros de hardware & Regla de preservación &
Comportamiento en subrutinas \\
\midrule
\endfirsthead
\toprule
Clase ABI & Registros de hardware & Regla de preservación &
Comportamiento en subrutinas \\
\midrule
\endhead
\bottomrule
\endfoot
\textbf{Callee-saved} (no volátiles) &
  \texttt{RBX}, \texttt{RBP}, \texttt{R12}--\texttt{R15}, \texttt{RSP} &
  Obligatorio preservar & Si la función los modifica, debe guardarlos en la
  pila y restaurarlos antes de retornar. \\
\textbf{Caller-saved} (volátiles) &
  \texttt{RAX}, \texttt{RCX}, \texttt{RDX}, \texttt{RSI}, \texttt{RDI},
  \texttt{R8}--\texttt{R11} & Libre modificación & La función llamada puede
  utilizarlos sin preservación; el invocador asume que se destruyen tras un
  \texttt{call}. \\
\textbf{SIMD volátiles} &
  \texttt{XMM0}--\texttt{XMM15} / \texttt{YMM0}--\texttt{YMM15} &
  Libre modificación & Todos los registros vectoriales son volátiles bajo la
  ABI System V AMD64. \\
\end{longtable}

\begin{longtable}{@{}>{\raggedright\arraybackslash}p{0.12\textwidth}
  >{\raggedright\arraybackslash}p{0.13\textwidth}
  >{\raggedright\arraybackslash}p{0.13\textwidth}
  >{\raggedright\arraybackslash}p{0.16\textwidth}
  >{\raggedright\arraybackslash}p{0.40\textwidth}@{}}
\caption{Asignación de registros de \texttt{sum\_array}: implementación escalar.}
\label{tab:sum-array-scalar-registers}\\
\toprule
Registro & Tipo y ancho & Estado ABI & Fase & Rol técnico y semántica \\
\midrule
\endfirsthead
\toprule
Registro & Tipo y ancho & Estado ABI & Fase & Rol técnico y semántica \\
\midrule
\endhead
\bottomrule
\endfoot
\texttt{RDI} & GPR 64 bits & Caller-saved & Entrada/bucle & Puntero base
  al arreglo de entrada \texttt{arr}. \\
\texttt{ESI} & GPR 32 bits & Caller-saved & Entrada/bucle & Límite superior
  $n$ y condición de parada. \\
\texttt{EAX}/\texttt{RAX} & GPR 32/64 bits & Caller-saved & Bucle escalar &
  Índice $i$; direccionamiento \texttt{[rdi + rax*4]}. \\
\texttt{XMM0} & SSE 128 bits & Caller-saved & Bucle/retorno & Acumulador
  escalar inicializado con \texttt{xorps}; valor de retorno en el carril 0. \\
\texttt{XMM1} & SSE 128 bits & Caller-saved & Bucle escalar & Carga temporal
  del elemento \texttt{arr[i]} mediante \texttt{movss}. \\
\end{longtable}

\begin{longtable}{@{}>{\raggedright\arraybackslash}p{0.12\textwidth}
  >{\raggedright\arraybackslash}p{0.13\textwidth}
  >{\raggedright\arraybackslash}p{0.13\textwidth}
  >{\raggedright\arraybackslash}p{0.17\textwidth}
  >{\raggedright\arraybackslash}p{0.39\textwidth}@{}}
\caption{Asignación de registros de \texttt{sum\_array}: implementación AVX2.}
\label{tab:sum-array-avx2-registers}\\
\toprule
Registro & Tipo y ancho & Estado ABI & Fase & Rol técnico y semántica \\
\midrule
\endfirsthead
\toprule
Registro & Tipo y ancho & Estado ABI & Fase & Rol técnico y semántica \\
\midrule
\endhead
\bottomrule
\endfoot
\texttt{RDI} & GPR 64 bits & Caller-saved & Entrada/bucles & Puntero base
  \texttt{arr}, alineado a 32 bytes. \\
\texttt{ESI} & GPR 32 bits & Caller-saved & Entrada/remanente & Cantidad total
  de elementos $n$. \\
\texttt{ECX} & GPR 32 bits & Caller-saved & Particionado/AVX2 & Límite
  $N_{\mathrm{vec}}=n\mathbin{\&}(\mathord{\sim}7)$. \\
\texttt{EAX}/\texttt{RAX} & GPR 32/64 bits & Caller-saved & AVX2/remanente &
  Índice $i$; avanza ocho elementos en AVX2 y uno en el remanente. \\
\texttt{YMM0} & AVX 256 bits & Caller-saved & Bucle AVX2 & Acumulador de
  ocho sumas parciales, inicializado con \texttt{vxorps}. \\
\texttt{YMM1} & AVX 256 bits & Caller-saved & Bucle AVX2 & Carga alineada
  de ocho valores mediante \texttt{vmovaps}. \\
\texttt{XMM2} & SSE 128 bits & Caller-saved & Reducción & Mitad alta de
  \texttt{YMM0}, extraída con \texttt{vextractf128}. \\
\texttt{XMM0} & SSE 128 bits & Caller-saved & Reducción/remanente/retorno &
  Acumulador colapsado con \texttt{vhaddps}; recibe el remanente y retorna el
  resultado en el carril 0. \\
\texttt{XMM1} & SSE 128 bits & Caller-saved & Remanente escalar & Carga del
  valor restante mediante \texttt{vmovss}. \\
\end{longtable}

\begin{longtable}{@{}>{\raggedright\arraybackslash}p{0.12\textwidth}
  >{\raggedright\arraybackslash}p{0.13\textwidth}
  >{\raggedright\arraybackslash}p{0.13\textwidth}
  >{\raggedright\arraybackslash}p{0.17\textwidth}
  >{\raggedright\arraybackslash}p{0.39\textwidth}@{}}
\caption{Asignación de registros de \texttt{compute\_stats}: implementación escalar.}
\label{tab:compute-stats-scalar-registers}\\
\toprule
Registro & Tipo y ancho & Estado ABI & Fase & Rol técnico y semántica \\
\midrule
\endfirsthead
\toprule
Registro & Tipo y ancho & Estado ABI & Fase & Rol técnico y semántica \\
\midrule
\endhead
\bottomrule
\endfoot
\texttt{R12} & GPR 64 bits & Callee-saved & Pasadas 1 y 2 & Puntero base
  \texttt{arr}, recibido inicialmente en \texttt{RDI}. \\
\texttt{R13D} & GPR 32 bits & Callee-saved & Pasadas 1 y 2 & Tamaño $n$,
  recibido inicialmente en \texttt{ESI}. \\
\texttt{R14} & GPR 64 bits & Callee-saved & Fin de pasada 1 & Puntero
  \texttt{mean*}, recibido en \texttt{RDX}. \\
\texttt{R15} & GPR 64 bits & Callee-saved & Fin de pasada 2 & Puntero
  \texttt{var*}, recibido en \texttt{RCX}. \\
\texttt{RBX} & GPR 64 bits & Callee-saved & Fin de pasada 1 & Puntero
  \texttt{min*}, recibido en \texttt{R8}. \\
\texttt{RBP} & GPR 64 bits & Callee-saved & Fin de pasada 1 & Puntero
  \texttt{max*}, recibido en \texttt{R9}. \\
\texttt{EAX}/\texttt{RAX} & GPR 32/64 bits & Caller-saved & Pasadas 1 y 2 &
  Índice $i$: $1\ldots n-1$ en P1 y $0\ldots n-1$ en P2. \\
\texttt{XMM0} & SSE 128 bits & Caller-saved & Pasada 1 & Acumulador de
  $\sum x_i$. \\
\texttt{XMM1} & SSE 128 bits & Caller-saved & Pasada 1 & Acumulador del
  mínimo mediante \texttt{minss}. \\
\texttt{XMM2} & SSE 128 bits & Caller-saved & Pasada 1 & Acumulador del
  máximo mediante \texttt{maxss}. \\
\texttt{XMM3} & SSE 128 bits & Caller-saved & Pasadas 1 y 2 & Registro
  temporal para \texttt{arr[i]} y la desviación respecto a la media. \\
\texttt{XMM4} & SSE 128 bits & Caller-saved & Pasadas 1 y 2 & Conversión de
  $n$ a punto flotante mediante \texttt{cvtsi2ss}. \\
\texttt{XMM5} & SSE 128 bits & Caller-saved & Fin P1/P2 & Media
  $\mu=\mathrm{sum}/n$, conservada para la segunda pasada. \\
\texttt{XMM6} & SSE 128 bits & Caller-saved & Pasada 2 & Acumulador de
  $\sum(x_i-\mu)^2$. \\
\end{longtable}

\begin{longtable}{@{}>{\raggedright\arraybackslash}p{0.12\textwidth}
  >{\raggedright\arraybackslash}p{0.13\textwidth}
  >{\raggedright\arraybackslash}p{0.13\textwidth}
  >{\raggedright\arraybackslash}p{0.17\textwidth}
  >{\raggedright\arraybackslash}p{0.39\textwidth}@{}}
\caption{Asignación de registros de \texttt{compute\_stats}: implementación AVX2.}
\label{tab:compute-stats-avx2-registers}\\
\toprule
Registro & Tipo y ancho & Estado ABI & Fase & Rol técnico y semántica \\
\midrule
\endfirsthead
\toprule
Registro & Tipo y ancho & Estado ABI & Fase & Rol técnico y semántica \\
\midrule
\endhead
\bottomrule
\endfoot
\texttt{R12} & GPR 64 bits & Callee-saved & Pasadas 1 y 2 & Puntero base
  \texttt{arr}, alineado a 32 bytes. \\
\texttt{R13D} & GPR 32 bits & Callee-saved & Pasadas 1 y 2 & Tamaño total $n$. \\
\texttt{R14} & GPR 64 bits & Callee-saved & Fin de pasada 1 & Destino
  \texttt{mean*}. \\
\texttt{R15} & GPR 64 bits & Callee-saved & Fin de pasada 2 & Destino
  \texttt{var*}. \\
\texttt{RBX} & GPR 64 bits & Callee-saved & Fin de pasada 1 & Destino
  \texttt{min*}. \\
\texttt{RBP} & GPR 64 bits & Callee-saved & Fin de pasada 1 & Destino
  \texttt{max*}. \\
\texttt{ECX} & GPR 32 bits & Caller-saved & Particionado P1/P2 & Límite
  $N_{\mathrm{vec}}=n\mathbin{\&}(\mathord{\sim}7)$. \\
\texttt{EAX}/\texttt{RAX} & GPR 32/64 bits & Caller-saved & Índice P1/P2 &
  Avanza ocho elementos en AVX2 y uno en el remanente. \\
\texttt{YMM0} & AVX 256 bits & Caller-saved & P1/P2 vectoriales & P1:
  acumulador de suma; P2: acumulador de desviaciones cuadráticas. \\
\texttt{YMM1} & AVX 256 bits & Caller-saved & P1/P2 vectoriales & P1:
  acumulador mínimo; P2: carga y cálculo de $(x-\mu)^2$. \\
\texttt{YMM2} & AVX 256 bits & Caller-saved & P1 vectorial & Acumulador máximo
  mediante \texttt{vmaxps}. \\
\texttt{YMM3} & AVX 256 bits & Caller-saved & P1/P2 vectoriales & P1: carga
  de ocho valores; P2: difusión de la media. \\
\texttt{XMM0} & SSE 128 bits & Caller-saved & Reducciones P1/P2 & Reducción
  de suma y de varianza. \\
\texttt{XMM1} & SSE 128 bits & Caller-saved & Reducción/remanente & Mínimo
  colapsado en P1 y temporal de diferencias en el remanente de P2. \\
\texttt{XMM2} & SSE 128 bits & Caller-saved & Reducción/remanente & Máximo
  colapsado y auxiliar escalar. \\
\texttt{XMM3} & SSE 128 bits & Caller-saved & Reducción/remanente & Auxiliar
  para \texttt{vshufps} y operaciones del remanente. \\
\texttt{XMM4} & SSE 128 bits & Caller-saved & Divisor P1/P2 & Representación
  flotante de $n$ mediante \texttt{vcvtsi2ss}. \\
\texttt{XMM5} & SSE 128 bits & Caller-saved & Pasadas 1 y 2 & Media escalar
  $\mu$, utilizada en la difusión y el remanente. \\
\end{longtable}

\begin{longtable}{@{}>{\raggedright\arraybackslash}p{0.12\textwidth}
  >{\raggedright\arraybackslash}p{0.13\textwidth}
  >{\raggedright\arraybackslash}p{0.13\textwidth}
  >{\raggedright\arraybackslash}p{0.17\textwidth}
  >{\raggedright\arraybackslash}p{0.39\textwidth}@{}}
\caption{Asignación de registros de \texttt{normalize\_array}: implementación escalar.}
\label{tab:normalize-scalar-registers}\\
\toprule
Registro & Tipo y ancho & Estado ABI & Fase & Rol técnico y semántica \\
\midrule
\endfirsthead
\toprule
Registro & Tipo y ancho & Estado ABI & Fase & Rol técnico y semántica \\
\midrule
\endhead
\bottomrule
\endfoot
\texttt{RDI} & GPR 64 bits & Caller-saved & Entrada/bucle & Puntero al arreglo
  de entrada \texttt{in}. \\
\texttt{RSI} & GPR 64 bits & Caller-saved & Entrada/bucle & Puntero al arreglo
  de salida \texttt{out}. \\
\texttt{EDX} & GPR 32 bits & Caller-saved & Entrada/bucle & Tamaño $n$. \\
\texttt{XMM0} & SSE 128 bits & Caller-saved & Entrada/bucle & Media escalar
  $\mu$. \\
\texttt{XMM1} & SSE 128 bits & Caller-saved & Entrada/bucle & Desviación estándar
  $\sigma$. \\
\texttt{XMM2} & SSE 128 bits & Caller-saved & Detección de borde & Cero escalar para
  \texttt{ucomiss xmm1, xmm2}. \\
\texttt{EAX}/\texttt{RAX} & GPR 32/64 bits & Caller-saved & Bucle & Índice
  $i=0\ldots n-1$. \\
\texttt{XMM3} & SSE 128 bits & Caller-saved & Bucle & Temporal que contiene
  \texttt{in[i]}, $\mathrm{in}[i]-\mu$ y el resultado normalizado. \\
\end{longtable}

\begin{longtable}{@{}>{\raggedright\arraybackslash}p{0.12\textwidth}
  >{\raggedright\arraybackslash}p{0.13\textwidth}
  >{\raggedright\arraybackslash}p{0.13\textwidth}
  >{\raggedright\arraybackslash}p{0.17\textwidth}
  >{\raggedright\arraybackslash}p{0.39\textwidth}@{}}
\caption{Asignación de registros de \texttt{normalize\_array}: implementación AVX2.}
\label{tab:normalize-avx2-registers}\\
\toprule
Registro & Tipo y ancho & Estado ABI & Fase & Rol técnico y semántica \\
\midrule
\endfirsthead
\toprule
Registro & Tipo y ancho & Estado ABI & Fase & Rol técnico y semántica \\
\midrule
\endhead
\bottomrule
\endfoot
\texttt{RDI} & GPR 64 bits & Caller-saved & Entrada/bucles & Puntero base
  \texttt{in}, alineado a 32 bytes. \\
\texttt{RSI} & GPR 64 bits & Caller-saved & Entrada/bucles & Puntero base
  \texttt{out}, alineado a 32 bytes. \\
\texttt{EDX} & GPR 32 bits & Caller-saved & Entrada/bucles & Tamaño $n$. \\
\texttt{XMM0} & SSE 128 bits & Caller-saved & Entrada/difusión & Media escalar
  $\mu$ recibida desde C. \\
\texttt{XMM1} & SSE 128 bits & Caller-saved & Entrada/difusión & Desviación estándar
  $\sigma$ recibida desde C. \\
\texttt{XMM2} & SSE 128 bits & Caller-saved & Detección/remanente & Comparación
  $\sigma=0.0$ y temporal escalar. \\
\texttt{YMM4} & AVX 256 bits & Caller-saved & AVX2/remanente & Broadcast
  $[\mu,\mu,\mu,\mu,\mu,\mu,\mu,\mu]$. \\
\texttt{YMM5} & AVX 256 bits & Caller-saved & AVX2/remanente & Broadcast
  $[\sigma,\sigma,\sigma,\sigma,\sigma,\sigma,\sigma,\sigma]$. \\
\texttt{ECX} & GPR 32 bits & Caller-saved & Particionado AVX2 & Límite
  $N_{\mathrm{vec}}=n\mathbin{\&}(\mathord{\sim}7)$. \\
\texttt{EAX}/\texttt{RAX} & GPR 32/64 bits & Caller-saved & AVX2/remanente &
  Índice $i$; paso de ocho en vector y uno en el remanente. \\
\texttt{YMM0} & AVX 256 bits & Caller-saved & Bucle AVX2 & Carga de ocho valores,
  resta, división y almacenamiento empaquetados. \\
\end{longtable}

\endgroup

\end{document}
